\input{preamble}

% OK, start here.
%
\begin{document}

\title{Critères de représentabilité}

\maketitle

\phantomsection
\label{section-phantom}

\tableofcontents




\section{Introduction}
\label{section-introduction}

\noindent
Le but de ce chapitre est de trouver des critères garantissant qu'un
champ en groupoïdes sur la catégorie des schémas munie de la topologie fppf
est un champ algébrique. Historiquement, cela consistait souvent à démontrer que
certains foncteurs étaient représentables ; voir les exposés de Grothendieck
\cite{Gr-I},
\cite{Gr-II},
\cite{Gr-III},
\cite{Gr-IV},
\cite{Gr-V}, and
\cite{Gr-VI}.
Cela explique le titre de ce chapitre. Les travaux d'Artin constituent une autre
source importante de ces résultats ; voir
\cite{ArtinI},
\cite{ArtinII},
\cite{Artin-Theorem-Representability},
\cite{Artin-Construction-Techniques},
\cite{Artin-Algebraic-Spaces},
\cite{Artin-Algebraic-Approximation},
\cite{Artin-Implicit-Function}, and
\cite{ArtinVersal}.

\medskip\noindent
Certaines notations, conventions et certains termes de ce chapitre sont peu
commodes et peuvent sembler inversés au lecteur expérimenté. C'est intentionnel.
On trouvera une explication dans Quot, section
\ref{quot-section-conventions}.



\section{Conventions}
\label{section-conventions}

\noindent
Les conventions employées dans ce chapitre sont les mêmes que dans le
chapitre consacré aux champs algébriques ; voir
Algebraic Stacks, section \ref{algebraic-section-conventions}.




\section{Ce que nous savons déjà}
\label{section-done-so-far}

\noindent
L'analogue de ce chapitre pour les espaces algébriques est le chapitre intitulé
« Bootstrap » ; voir
Bootstrap, section \ref{bootstrap-section-introduction}.
Ce chapitre contient déjà certains résultats de représentabilité.
En outre, nous avons déjà établi dans le chapitre sur les champs algébriques
une partie des résultats préliminaires qui y sont traités.
En voici la liste :
\begin{enumerate}
\item Les morphismes de préfaisceaux représentables par des espaces algébriques sont étudiés dans
Bootstrap, section
\ref{bootstrap-section-morphism-representable-by-spaces}.
Dans
Algebraic Stacks, section
\ref{algebraic-section-morphisms-representable-by-algebraic-spaces}
nous étudions la notion de $1$-morphisme de catégories fibrées en groupoïdes
représentable par des espaces algébriques.
\item Les propriétés des morphismes de préfaisceaux représentables par des
espaces algébriques sont étudiées dans
Bootstrap, section
\ref{bootstrap-section-representable-by-spaces-properties}.
Dans
Algebraic Stacks, section
\ref{algebraic-section-representable-properties}
nous étudions les propriétés des $1$-morphismes de catégories fibrées en groupoïdes
représentables par des espaces algébriques.
\item Nous avons démontré que, si $F$ est un faisceau dont la diagonale est représentable
par des espaces algébriques et qui possède un recouvrement \'etale par un espace
algébrique, alors $F$ est un espace algébrique ; voir
Bootstrap, théorème \ref{bootstrap-theorem-bootstrap}.
(Il s'agit d'une version faible du résultat de l'alinéa suivant.)
\item
\label{item-bootstrap-final}
Nous avons démontré que, si $F$ est un faisceau et s'il existe un espace
algébrique $U$ et un morphisme $U \to F$ représentable par des espaces
algébriques, surjectif, plat et localement de présentation finie, alors
$F$ est un espace algébrique ; voir
Bootstrap, théorème \ref{bootstrap-theorem-final-bootstrap}.
\item Nous avons aussi démontré l'analogue « lisse » de
(\ref{item-bootstrap-final}) pour les champs
algébriques : si $\mathcal{X}$ est un champ en groupoïdes sur
$(\Sch/S)_{fppf}$ et s'il existe un champ en groupoïdes
$\mathcal{U}$ sur $(\Sch/S)_{fppf}$ représentable
par un espace algébrique ainsi qu'un $1$-morphisme $u : \mathcal{U} \to \mathcal{X}$
représentable par des espaces algébriques, surjectif et lisse,
alors $\mathcal{X}$ est un champ algébrique ; voir
Algebraic Stacks, lemme
\ref{algebraic-lemma-smooth-surjective-morphism-implies-algebraic}.
\end{enumerate}
Notre première tâche consiste maintenant à démontrer l'analogue de
(\ref{item-bootstrap-final}) pour les champs
algébriques en général ; c'est le
théorème \ref{theorem-bootstrap}.



\section{Morphismes de champs en groupoïdes}
\label{section-1-morphisms}

\noindent
Cette section est préliminaire et peut être omise en première lecture.

\begin{lemma}
\label{lemma-etale-permanence}
Soient $\mathcal{X} \to \mathcal{Y} \to \mathcal{Z}$ des
$1$-morphismes de catégories fibrées en groupoïdes sur
$(\Sch/S)_{fppf}$.
Si $\mathcal{X} \to \mathcal{Z}$ et $\mathcal{Y} \to \mathcal{Z}$ sont
représentables par des espaces algébriques et étales, il en est de même de
$\mathcal{X} \to \mathcal{Y}$.
\end{lemma}

\begin{proof}
Soit $\mathcal{U}$ une catégorie fibrée en groupoïdes représentable sur $S$.
Soit $f : \mathcal{U} \to \mathcal{Y}$ un $1$-morphisme. Il faut montrer que
$\mathcal{X} \times_\mathcal{Y} \mathcal{U}$ est représentable par un
espace algébrique et étale sur $\mathcal{U}$.
Considérons la composée $h : \mathcal{U} \to \mathcal{Z}$. Alors
$$
\mathcal{X} \times_\mathcal{Z} \mathcal{U}
\longrightarrow
\mathcal{Y} \times_\mathcal{Z} \mathcal{U}
$$
est un $1$-morphisme entre deux catégories fibrées en groupoïdes, toutes deux
représentables par des espaces algébriques et étales sur $\mathcal{U}$.
Ainsi, d'après
Properties of Spaces, lemme \ref{spaces-properties-lemma-etale-permanence},
il est représenté par un morphisme étale d'espaces algébriques.
Enfin, on obtient le résultat voulu puisque le morphisme $f$ induit
un morphisme $\mathcal{U} \to \mathcal{Y} \times_\mathcal{Z} \mathcal{U}$
et que l'on a
$$
\mathcal{X} \times_\mathcal{Y} \mathcal{U} =
(\mathcal{X} \times_\mathcal{Z} \mathcal{U})
\times_{(\mathcal{Y} \times_\mathcal{Z} \mathcal{U})}
\mathcal{U}.
$$
\end{proof}

\begin{lemma}
\label{lemma-stack-in-setoids-descent}
Soient $\mathcal{X}, \mathcal{Y}, \mathcal{Z}$ des champs en groupoïdes
sur $(\Sch/S)_{fppf}$. Supposons que $\mathcal{X} \to \mathcal{Z}$
et $\mathcal{Y} \to \mathcal{Z}$ soient des $1$-morphismes.
Si
\begin{enumerate}
\item $\mathcal{Y}$ et $\mathcal{Z}$ sont représentables par des espaces algébriques
$Y$ et $Z$ sur $S$,
\item le morphisme associé d'espaces algébriques $Y \to Z$ est surjectif,
plat et localement de présentation finie, et
\item $\mathcal{Y} \times_\mathcal{Z} \mathcal{X}$ est un champ en
groupoïdes discrets,
\end{enumerate}
alors $\mathcal{X}$ est un champ en groupoïdes discrets.
\end{lemma}

\begin{proof}
C'est un cas particulier de
Stacks, lemme \ref{stacks-lemma-stack-in-setoids-descent}.
\end{proof}

\noindent
Le lemme suivant est l'analogue de
Algebraic Stacks, lemme
\ref{algebraic-lemma-smooth-surjective-morphism-implies-algebraic}
et sera supplanté par le
théorème \ref{theorem-bootstrap}, plus fort.

\begin{lemma}
\label{lemma-flat-finite-presentation-surjective-diagonal}
Soit $S$ un schéma.
Soit $u : \mathcal{U} \to \mathcal{X}$ un $1$-morphisme de
champs en groupoïdes sur $(\Sch/S)_{fppf}$. Si
\begin{enumerate}
\item $\mathcal{U}$ est représentable par un espace algébrique, et
\item $u$ est représentable par des espaces algébriques, surjectif, plat et
localement de présentation finie,
\end{enumerate}
alors
$\Delta : \mathcal{X} \to \mathcal{X} \times \mathcal{X}$
est représentable par des espaces algébriques.
\end{lemma}

\begin{proof}
Étant donnés deux schémas $T_1$, $T_2$ sur $S$, notons
$\mathcal{T}_i = (\Sch/T_i)_{fppf}$ les catégories fibrées représentables
associées. Supposons donnés des $1$-morphismes
$f_i : \mathcal{T}_i \to \mathcal{X}$.
D'après
Algebraic Stacks, lemme \ref{algebraic-lemma-representable-diagonal},
il suffit de démontrer que le $2$-produit
fibré $\mathcal{T}_1 \times_\mathcal{X} \mathcal{T}_2$
est représentable par un espace algébrique. D'après
Stacks, lemme
\ref{stacks-lemma-2-fibre-product-stacks-in-setoids-over-stack-in-groupoids}
c'est en tout cas un champ en groupoïdes discrets. Ainsi,
$\mathcal{T}_1 \times_\mathcal{X} \mathcal{T}_2$ correspond
à un faisceau $F$ sur $(\Sch/S)_{fppf}$ ; voir
Stacks, lemme \ref{stacks-lemma-stack-in-setoids-characterize}.
Soit $U$ l'espace algébrique qui représente $\mathcal{U}$.
Par hypothèse,
$$
\mathcal{T}_i' = \mathcal{U} \times_{u, \mathcal{X}, f_i} \mathcal{T}_i
$$
est représentable par un espace algébrique $T'_i$ sur $S$. Par conséquent,
$\mathcal{T}_1' \times_\mathcal{U} \mathcal{T}_2'$ est représentable
par l'espace algébrique $T'_1 \times_U T'_2$.
Considérons le diagramme commutatif
$$
\xymatrix{
&
\mathcal{T}_1 \times_{\mathcal X} \mathcal{T}_2 \ar[rr]\ar'[d][dd] & &
\mathcal{T}_1 \ar[dd] \\
\mathcal{T}_1' \times_\mathcal{U} \mathcal{T}_2' \ar[ur]\ar[rr]\ar[dd] & &
\mathcal{T}_1' \ar[ur]\ar[dd] \\
&
\mathcal{T}_2 \ar'[r][rr] & &
\mathcal X \\
\mathcal{T}_2' \ar[rr]\ar[ur] & &
\mathcal{U} \ar[ur] }
$$
Dans ce diagramme, les carrés inférieur, droit, arrière et
avant sont des $2$-produits fibrés. Un argument formel montre alors
que $\mathcal{T}_1' \times_\mathcal{U} \mathcal{T}_2' \to
\mathcal{T}_1 \times_{\mathcal X} \mathcal{T}_2$
est le « changement de base » de $\mathcal{U} \to \mathcal{X}$ ; plus précisément,
le diagramme
$$
\xymatrix{
\mathcal{T}_1' \times_\mathcal{U} \mathcal{T}_2' \ar[d] \ar[r] &
\mathcal{U} \ar[d] \\
\mathcal{T}_1 \times_{\mathcal X} \mathcal{T}_2 \ar[r] &
\mathcal{X}
}
$$
est un carré $2$-cartésien.
Ainsi, $T'_1 \times_U T'_2 \to F$ est représentable par des espaces algébriques,
plat, localement de présentation finie et surjectif ; voir
Algebraic Stacks, lemmes
\ref{algebraic-lemma-map-fibred-setoids-representable-algebraic-spaces},
\ref{algebraic-lemma-base-change-representable-by-spaces},
\ref{algebraic-lemma-map-fibred-setoids-property}, and
\ref{algebraic-lemma-base-change-representable-transformations-property}.
Par conséquent, $F$ est un espace algébrique d'après
Bootstrap, théorème \ref{bootstrap-theorem-final-bootstrap},
ce qui conclut.
\end{proof}

\begin{lemma}
\label{lemma-second-diagonal}
Soit $\mathcal{X}$ une catégorie fibrée en groupoïdes sur $(\Sch/S)_{fppf}$.
Les assertions suivantes sont équivalentes :
\begin{enumerate}
\item $\Delta_\Delta : \mathcal{X} \to
\mathcal{X} \times_{\mathcal{X} \times \mathcal{X}} \mathcal{X}$
est représentable par des espaces algébriques,
\item pour tout $1$-morphisme $\mathcal{V} \to \mathcal{X} \times \mathcal{X}$
où $\mathcal{V}$ est représentable (par un schéma), le produit fibré
$\mathcal{Y} =
\mathcal{X} \times_{\Delta, \mathcal{X} \times \mathcal{X}} \mathcal{V}$
a une diagonale représentable par des espaces algébriques.
\end{enumerate}
\end{lemma}

\begin{proof}
Bien que cela demande un peu de réflexion, l'argument est entièrement formel.
Rappelons en effet que
$\mathcal{X} \times_{\mathcal{X} \times \mathcal{X}} \mathcal{X} =
\mathcal{I}_\mathcal{X}$ est l'inertie de $\mathcal{X}$ et que
$\Delta_\Delta$ est la section unité de $\mathcal{I}_\mathcal{X}$ ; voir
Categories, section \ref{categories-section-inertia}.
La condition (1) signifie donc ceci : étant donnés un schéma $V$, un objet $x$ de
$\mathcal{X}$ au-dessus de $V$ et un morphisme $\alpha : x \to x$ de $\mathcal{X}_V$,
la condition « $\alpha = \text{id}_x$ » définit un espace algébrique sur $V$.
(Autrement dit, il existe un monomorphisme d'espaces algébriques $W \to V$
tel qu'un morphisme de schémas $f : T \to V$ se factorise par $W$
si et seulement si $f^*\alpha = \text{id}_{f^*x}$.)

\medskip\noindent
Soient d'autre part $V$ un schéma et $x, y$ des objets de
$\mathcal{X}$ au-dessus de $V$. Alors $(x, y)$ définit un morphisme
$\mathcal{V} = (\Sch/V)_{fppf} \to \mathcal{X} \times \mathcal{X}$.
Soient ensuite $h : V' \to V$ un morphisme de schémas et
$\alpha : h^*x \to h^*y$ et $\beta : h^*x \to h^*y$ des morphismes
de $\mathcal{X}_{V'}$. Alors $(\alpha, \beta)$ définit un morphisme
$\mathcal{V}' = (\Sch/V')_{fppf} \to \mathcal{Y} \times \mathcal{Y}$.
La condition (2) affirme alors que, quels que soient les choix ci-dessus, la
condition « $\alpha = \beta$ » définit un espace algébrique sur $V'$.

\medskip\noindent
Pour voir l'équivalence, étant donné $(\alpha, \beta)$ comme en (2), on constate que
(1) implique que « $\alpha^{-1} \circ \beta = \text{id}_{h^*x}$ »
définit un espace algébrique. L'implication (2) $\Rightarrow$ (1)
s'obtient en prenant $h = \text{id}_V$ et $\beta = \text{id}_x$.
\end{proof}











\section{Commutation aux limites sur les objets}
\label{section-limit-preserving}

\noindent
Soit $S$ un schéma. Soit $p : \mathcal{X} \to \mathcal{Y}$ un $1$-morphisme
de catégories fibrées en groupoïdes sur $(\Sch/S)_{fppf}$. Nous dirons que
$p$ {\it commute aux limites sur les objets} si la condition suivante est satisfaite :
pour toute donnée formée de
\begin{enumerate}
\item un schéma affine $U = \lim_{i \in I} U_i$ écrit comme
limite projective filtrante de schémas affines $U_i$ sur $S$,
\item un objet $y_i$ de $\mathcal{Y}$ au-dessus de $U_i$ pour un certain $i$,
\item un objet $x$ de $\mathcal{X}$ au-dessus de $U$, et
\item un isomorphisme $\gamma : p(x) \to y_i|_U$,
\end{enumerate}
il existe un $i' \geq i$, un objet $x_{i'}$ de
$\mathcal{X}$ au-dessus de $U_{i'}$, un isomorphisme
$\beta : x_{i'}|_U \to x$ et un isomorphisme
$\gamma_{i'} : p(x_{i'}) \to y_i|_{U_{i'}}$
tels que
\begin{equation}
\label{equation-limit-preserving}
\vcenter{
\xymatrix{
p(x_{i'}|_U) \ar[d]_{p(\beta)} \ar[rr]_{\gamma_{i'}|_U} & &
(y_i|_{U_{i'}})|_U \ar@{=}[d] \\
p(x) \ar[rr]^\gamma & & y_i|_U
}
}
\end{equation}
soit commutatif. Dans cette situation, on dit que « $(i', x_{i'}, \beta, \gamma_{i'})$
est une {\it solution} du problème posé par les données (1), (2), (3), (4) ».
La motivation de cette définition vient de
Limits of Spaces,
lemme \ref{spaces-limits-lemma-characterize-relative-limit-preserving}.

\begin{lemma}
\label{lemma-base-change-limit-preserving}
Soient $p : \mathcal{X} \to \mathcal{Y}$ et $q : \mathcal{Z} \to \mathcal{Y}$
des $1$-morphismes de catégories fibrées en groupoïdes sur $(\Sch/S)_{fppf}$.
Si $p : \mathcal{X} \to \mathcal{Y}$ commute aux limites sur les objets, il en
est de même du changement de base
$p' : \mathcal{X} \times_\mathcal{Y} \mathcal{Z} \to \mathcal{Z}$
de $p$ par $q$.
\end{lemma}

\begin{proof}
L'assertion est formelle. Soit $U = \lim_{i \in I} U_i$ la limite projective filtrante
de schémas affines $U_i$ sur $S$, soit $z_i$ un objet de $\mathcal{Z}$
au-dessus de $U_i$ pour un certain $i$, soit $w$ un objet de
$\mathcal{X} \times_\mathcal{Y} \mathcal{Z}$ au-dessus de $U$, et soit
$\delta : p'(w) \to z_i|_U$ un isomorphisme.
On peut écrire
$w = (U, x, z, \alpha)$, où $x$ est un objet de $\mathcal{X}$ au-dessus de $U$,
$z$ un objet de $\mathcal{Z}$ au-dessus de $U$ et
$\alpha : p(x) \to q(z)$ un isomorphisme. Puisque $p'(w) = z$, on a
$\delta : z \to z_i|_U$. Posons $y_i = q(z_i)$ et
$\gamma = q(\delta) \circ \alpha : p(x) \to y_i|_U$.
Comme $p$ commute aux limites sur les objets, il existe $i' \geq i$
et un objet $x_{i'}$ de $\mathcal{X}$ au-dessus de $U_{i'}$, ainsi que des
isomorphismes $\beta : x_{i'}|_U \to x$ et
$\gamma_{i'} : p(x_{i'}) \to y_i|_{U_{i'}}$ tels que
(\ref{equation-limit-preserving}) soit commutatif. Considérons alors l'objet
$w_{i'} = (U_{i'}, x_{i'}, z_i|_{U_{i'}}, \gamma_{i'})$ de
$\mathcal{X} \times_\mathcal{Y} \mathcal{Z}$ au-dessus de $U_{i'}$
et définissons les isomorphismes
$$
w_{i'}|_U = (U, x_{i'}|_U, z_i|_U, \gamma_{i'}|_U)
\xrightarrow{(\beta, \delta^{-1})}
(U, x, z, \alpha) = w
$$
et
$$
p'(w_{i'}) = z_i|_{U_{i'}} \xrightarrow{\text{id}} z_i|_{U_{i'}}.
$$
Ceux-ci se combinent en une solution du problème.
\end{proof}

\begin{lemma}
\label{lemma-composition-limit-preserving}
Soient $p : \mathcal{X} \to \mathcal{Y}$ et $q : \mathcal{Y} \to \mathcal{Z}$
des $1$-morphismes de catégories fibrées en groupoïdes sur $(\Sch/S)_{fppf}$.
Si $p$ et $q$ commutent aux limites sur les objets, il en est de même de la composée
$q \circ p$.
\end{lemma}

\begin{proof}
L'assertion est formelle. Soit $U = \lim_{i \in I} U_i$ la limite projective filtrante
de schémas affines $U_i$ sur $S$, soit $z_i$ un objet de $\mathcal{Z}$
au-dessus de $U_i$ pour un certain $i$, soit $x$ un objet de $\mathcal{X}$ au-dessus de $U$,
et soit $\gamma : q(p(x)) \to z_i|_U$ un isomorphisme. Comme $q$
commute aux limites sur les objets, il existe $i' \geq i$, un objet
$y_{i'}$ de $\mathcal{Y}$ au-dessus de $U_{i'}$, un isomorphisme
$\beta : y_{i'}|_U \to p(x)$ et un isomorphisme
$\gamma_{i'} : q(y_{i'}) \to z_i|_{U_{i'}}$
tels que (\ref{equation-limit-preserving}) soit commutatif. Comme $p$
commute aux limites sur les objets, il existe $i'' \geq i'$, un objet
$x_{i''}$ de $\mathcal{X}$ au-dessus de $U_{i''}$, un isomorphisme
$\beta' : x_{i''}|_U \to x$ et un isomorphisme
$\gamma'_{i''} : p(x_{i''}) \to y_{i'}|_{U_{i''}}$
tels que (\ref{equation-limit-preserving}) soit commutatif.
La solution consiste à prendre $x_{i''}$ au-dessus de $U_{i''}$ avec l'isomorphisme
$$
q(p(x_{i''})) \xrightarrow{q(\gamma'_{i''})}
q(y_{i'})|_{U_{i''}} \xrightarrow{\gamma_{i'}|_{U_{i''}}}
z_i|_{U_{i''}}
$$
et l'isomorphisme $\beta' : x_{i''}|_U \to x$. Nous omettons de vérifier
que (\ref{equation-limit-preserving}) est commutatif.
\end{proof}

\begin{lemma}
\label{lemma-representable-by-spaces-limit-preserving}
Soit $p : \mathcal{X} \to \mathcal{Y}$ un $1$-morphisme de catégories
fibrées en groupoïdes sur $(\Sch/S)_{fppf}$. Si $p$ est
représentable par des espaces algébriques, les assertions suivantes sont équivalentes :
\begin{enumerate}
\item $p$ commute aux limites sur les objets, et
\item $p$ est localement de présentation finie (voir
Algebraic Stacks,
définition \ref{algebraic-definition-relative-representable-property}).
\end{enumerate}
\end{lemma}

\begin{proof}
Supposons (2). Soit $U = \lim_{i \in I} U_i$ la limite projective filtrante
de schémas affines $U_i$ sur $S$, soit $y_i$ un objet de $\mathcal{Y}$
au-dessus de $U_i$ pour un certain $i$, soit $x$ un objet de $\mathcal{X}$ au-dessus de $U$,
et soit $\gamma : p(x) \to y_i|_U$ un isomorphisme. Notons
$X_{y_i}$ un espace algébrique sur $U_i$ représentant le $2$-produit
fibré
$$
(\Sch/U_i)_{fppf} \times_{y_i, \mathcal{Y}, p} \mathcal{X}.
$$
Notons que $\xi = (U, U \to U_i, x, \gamma^{-1})$ définit un objet de
ce $2$-produit fibré au-dessus de $U$. Par le lemme de $2$-Yoneda, $\xi$ correspond
à un morphisme $f_\xi : U \to X_{y_i}$ sur $U_i$. D'après
Limits of Spaces, Proposition
\ref{spaces-limits-proposition-characterize-locally-finite-presentation}
il existe $i' \geq i$ et un morphisme $f_{i'} : U_{i'} \to X_{y_i}$
tels que $f_\xi$ soit la composée de $f_{i'}$ et du morphisme de
projection $U \to U_{i'}$. En outre, le lemme de $2$-Yoneda montre que
$f_{i'}$ correspond à un objet
$\xi_{i'} = (U_{i'}, U_{i'} \to U_i, x_{i'}, \alpha)$ of
du $2$-produit fibré affiché au-dessus de $U_{i'}$, dont la restriction à
$U$ redonne $\xi$. En particulier, on obtient un isomorphisme
$\beta : x_{i'}|U \to x$. Notons que $\alpha : y_i|_{U_{i'}} \to p(x_{i'})$.
Ainsi, prendre $x_{i'}$, l'isomorphisme
$\beta : x_{i'}|U \to x$ et l'isomorphisme
$\gamma_{i'} = \alpha^{-1} : p(x_{i'}) \to y_i|_{U_{i'}}$
fournit une solution du problème.

\medskip\noindent
Supposons (1). Choisissons un schéma $T$ et un $1$-morphisme
$y : (\Sch/T)_{fppf} \to \mathcal{Y}$. Soit
$X_y$ un espace algébrique sur $T$ représentant le $2$-produit fibré
$(\Sch/T)_{fppf} \times_{y, \mathcal{Y}, p} \mathcal{X}$.
Il faut montrer que $X_y \to T$ est localement de présentation finie.
Pour cela, nous utiliserons le critère de
Limits of Spaces, remarque \ref{spaces-limits-remark-limit-preserving}.
Considérons un schéma affine $U = \lim_{i \in I} U_i$ écrit comme
limite projective filtrante de schémas affines sur $T$.
Choisissons $i \in I$ et posons $y_i = y|_{U_i}$. Notons aussi $i'$ un élément
de $I$ supérieur ou égal à $i$. D'après le lemme de $2$-Yoneda, les
morphismes $U \to X_y$ sur $T$ correspondent bijectivement
aux classes d'isomorphisme de couples $(x, \alpha)$, où $x$ est un objet
de $\mathcal{X}$ au-dessus de $U$ et $\alpha : y|_U \to p(x)$ un isomorphisme.
Bien entendu, donner $\alpha$ revient, à inversion près, à donner
un isomorphisme $\gamma : p(x) \to y_i|_U$.
Il en va de même pour les morphismes $U_{i'} \to X_y$ sur $T$. Ainsi, (1) garantit
que l'application canonique
$$
\colim_{i' \geq i} X_y(U_{i'}) \longrightarrow X_y(U)
$$
est surjective dans cette situation. Il résulte de
Limits of Spaces, lemme \ref{spaces-limits-lemma-surjection-is-enough},
que $X_y \to T$ est localement de présentation finie.
\end{proof}

\begin{lemma}
\label{lemma-open-immersion-limit-preserving}
Soit $p : \mathcal{X} \to \mathcal{Y}$ un $1$-morphisme de catégories
fibrées en groupoïdes sur $(\Sch/S)_{fppf}$. Supposons $p$ représentable
par des espaces algébriques et une immersion ouverte. Alors $p$ commute aux
limites sur les objets.
\end{lemma}

\begin{proof}
Cela résulte du
lemme \ref{lemma-representable-by-spaces-limit-preserving}
et, par le principe général de
Algebraic Stacks, lemme
\ref{algebraic-lemma-representable-transformations-property-implication})
du fait qu'une immersion ouverte d'espaces algébriques est
localement de présentation finie ; voir
Morphisms of Spaces, lemme
\ref{spaces-morphisms-lemma-open-immersion-locally-finite-presentation}.
\end{proof}

\noindent
Soit $S$ un schéma. Dans le lemme suivant, nous aurons besoin de la notion de
{\it taille} d'un espace algébrique $X$ sur $S$. Étant donné un cardinal
$\kappa$, nous dirons que $X$ vérifie $\text{size}(X) \leq \kappa$ si et seulement
s'il existe un schéma $U$ tel que $\text{size}(U) \leq \kappa$ (voir
Sets, section \ref{sets-section-categories-schemes}) et un morphisme
étale surjectif $U \to X$.

\begin{lemma}
\label{lemma-check-representable-limit-preserving}
Soit $S$ un schéma.
Soit $\kappa = \text{size}(T)$ pour un certain $T \in \Ob((\Sch/S)_{fppf})$.
Soit $f : \mathcal{X} \to \mathcal{Y}$ un $1$-morphisme
de catégories fibrées en groupoïdes sur $(\Sch/S)_{fppf}$
tel que
\begin{enumerate}
\item $\mathcal{Y} \to (\Sch/S)_{fppf}$ commute aux limites sur les objets,
\item pour tout schéma affine $V$ localement de présentation finie sur $S$ et
tout $y \in \Ob(\mathcal{Y}_V)$, le produit fibré
$(\Sch/V)_{fppf} \times_{y, \mathcal{Y}} \mathcal{X}$ est représentable
par un espace algébrique de taille $\leq \kappa$\footnote{Ceux qui négligent les
questions ensemblistes peuvent supprimer la condition sur la taille.},
\item $\mathcal{X}$ et $\mathcal{Y}$ sont des champs pour la topologie de Zariski.
\end{enumerate}
Alors $f$ est représentable par des espaces algébriques.
\end{lemma}

\begin{proof}
Soient $V$ un schéma sur $S$ et $y \in \mathcal{Y}_V$. Il faut démontrer que
$(\Sch/V)_{fppf} \times_{y, \mathcal{Y}} \mathcal{X}$ est représentable
par un espace algébrique.

\medskip\noindent
Cas I : $V$ est affine et s'envoie dans un ouvert affine $\Spec(\Lambda) \subset S$.
On peut alors écrire $V = \lim V_i$, chaque $V_i$ étant affine et de
présentation finie sur $\Spec(\Lambda)$ ; voir
Algebra, lemme \ref{algebra-lemma-ring-colimit-fp}.
L'hypothèse (1) montre alors que $y$ provient d'un objet $y_i$ au-dessus de $V_i$ pour un certain $i$.
Par l'hypothèse (3), le produit fibré
$(\Sch/V_i)_{fppf} \times_{y_i, \mathcal{Y}} \mathcal{X}$ est représentable
par un espace algébrique $Z_i$. Alors
$(\Sch/V)_{fppf} \times_{y, \mathcal{Y}} \mathcal{X}$ est représentable
par $Z_i \times_{V_i} V$.

\medskip\noindent
Cas II : $V$ est quelconque. Choisissons un recouvrement ouvert affine
$V = \bigcup_{i \in I} V_i$ tel que chaque $V_i$ s'envoie dans un ouvert affine
de $S$. Nous affirmons d'abord
que $\mathcal{Z} = (\Sch/V)_{fppf} \times_{y, \mathcal{Y}} \mathcal{X}$
est un champ en groupoïdes discrets pour la topologie de Zariski. En effet, c'est un champ en
groupoïdes pour la topologie de Zariski d'après
Stacks, lemme \ref{stacks-lemma-2-product-stacks-in-groupoids}.
Supposons ensuite que $z$ soit un objet de $\mathcal{Z}$ au-dessus d'un schéma $T$.
Notons $g : T \to V$ le morphisme correspondant à la
projection de $z$ dans $(\Sch/V)_{fppf}$. Considérons le faisceau de Zariski
$\mathit{I} = \mathit{Isom}_{\mathcal{Z}}(z, z)$. Le cas I montre que
$\mathit{I}|_{g^{-1}(V_i)} = *$ (le faisceau réduit à un élément). Ainsi,
$\mathit{I} = *$. Par conséquent, $\mathcal{Z}$ est fibrée en groupoïdes discrets. Pour achever
la démonstration, il faut montrer que le faisceau de Zariski
$Z : T \mapsto \Ob(\mathcal{Z}_T)/\cong$ est un espace algébrique ; voir
Algebraic Stacks, lemme
\ref{algebraic-lemma-characterize-representable-by-space}.
Il existe une application $p : Z \to V$ (une transformation de foncteurs) et le cas I
montre que $Z_i = p^{-1}(V_i)$ est un espace algébrique. Les morphismes
$Z_i \to Z$ sont représentables par des immersions ouvertes et
$\coprod Z_i \to Z$ est surjectif (pour la topologie de Zariski).
Ainsi, $Z$ est un faisceau pour la topologie fppf d'après
Bootstrap, lemme \ref{bootstrap-lemma-glueing-sheaves}.
Le lemme de Spaces, \ref{spaces-lemma-glueing-algebraic-spaces},
s'applique donc et montre que $Z$ est un espace algébrique\footnote{
Pour vérifier la condition ensembliste de ce lemme,
raisonnons comme suit. Choisissons d'abord le recouvrement ouvert de sorte que
$|I| \leq \text{size}(V)$. Choisissons ensuite des schémas $U_i$ de taille
$\leq \max(\kappa, \text{size}(V))$ et des morphismes étales surjectifs
$U_i \to Z_i$ ; cela est possible par l'hypothèse (2) et
Sets, lemme \ref{sets-lemma-bound-size-fibre-product}
(nous omettons les détails). Alors
Sets, lemme \ref{sets-lemma-what-is-in-it},
implique que $\coprod U_i$ est un objet de $(\Sch/S)_{fppf}$.
Par conséquent, $\coprod Z_i$ est un espace algébrique d'après
Spaces, lemme \ref{spaces-lemma-coproduct-algebraic-spaces}.
}.
\end{proof}

\begin{lemma}
\label{lemma-check-property-limit-preserving}
Soit $S$ un schéma. Soit $f : \mathcal{X} \to \mathcal{Y}$ un $1$-morphisme
de catégories fibrées en groupoïdes sur $(\Sch/S)_{fppf}$. Soit $\mathcal{P}$
une propriété des morphismes d'espaces algébriques comme dans
Algebraic Stacks, définition
\ref{algebraic-definition-relative-representable-property}. If
\begin{enumerate}
\item $f$ est représentable par des espaces algébriques,
\item $\mathcal{Y} \to (\Sch/S)_{fppf}$ commute aux limites sur les objets,
\item pour tout schéma affine $V$ localement de présentation finie sur $S$ et
tout $y \in \mathcal{Y}_V$, le morphisme d'espaces algébriques qui en résulte
$f_y : F_y \to V$, voir Algebraic Stacks, équation
(\ref{algebraic-equation-representable-by-algebraic-spaces}),
a la propriété $\mathcal{P}$.
\end{enumerate}
Alors $f$ a la propriété $\mathcal{P}$.
\end{lemma}

\begin{proof}
Soient $V$ un schéma sur $S$ et $y \in \mathcal{Y}_V$. Il faut montrer
que $F_y \to V$ a la propriété $\mathcal{P}$. Comme $\mathcal{P}$ est
locale pour la topologie fppf sur la base, on peut supposer que $V$ est un schéma affine qui
s'envoie dans un ouvert affine $\Spec(\Lambda) \subset S$. On peut donc écrire
$V = \lim V_i$, chaque $V_i$ étant affine et de présentation finie sur
$\Spec(\Lambda)$ ; voir Algebra, lemme \ref{algebra-lemma-ring-colimit-fp}.
L'hypothèse (2) montre alors que $y$ provient d'un objet $y_i$ au-dessus de $V_i$ pour un certain $i$.
Par l'hypothèse (3), le morphisme $F_{y_i} \to V_i$ a la propriété $\mathcal{P}$.
Comme $\mathcal{P}$ est stable par tout changement de base et que
$F_y = F_{y_i} \times_{V_i} V$, on conclut que $F_y \to V$ a la propriété
$\mathcal{P}$, comme voulu.
\end{proof}



\section{Formellement lisse sur les objets}
\label{section-formally-smooth}

\noindent
Soit $S$ un schéma. Soit $p : \mathcal{X} \to \mathcal{Y}$ un $1$-morphisme
de catégories fibrées en groupoïdes sur $(\Sch/S)_{fppf}$. Nous dirons que
$p$ est {\it formellement lisse sur les objets} si la condition suivante est satisfaite :
pour toute donnée formée de
\begin{enumerate}
\item un épaississement du premier ordre $U \subset U'$ de schémas affines sur $S$,
\item un objet $y'$ de $\mathcal{Y}$ au-dessus de $U'$,
\item un objet $x$ de $\mathcal{X}$ au-dessus de $U$, et
\item un isomorphisme $\gamma : p(x) \to y'|_U$,
\end{enumerate}
il existe un objet $x'$ de
$\mathcal{X}$ au-dessus de $U'$, muni d'un isomorphisme
$\beta : x'|_U \to x$ et d'un isomorphisme $\gamma' : p(x') \to y'$
tels que
\begin{equation}
\label{equation-formally-smooth}
\vcenter{
\xymatrix{
p(x'|_U) \ar[d]_{p(\beta)} \ar[rr]_{\gamma'|_U} & &
y'|_U \ar@{=}[d] \\
p(x) \ar[rr]^\gamma & & y'|_U
}
}
\end{equation}
soit commutatif. Dans cette situation, on dit que « $(x', \beta, \gamma')$
est une {\it solution} du problème posé par les données (1), (2), (3), (4) ».

\begin{lemma}
\label{lemma-base-change-formally-smooth}
Soient $p : \mathcal{X} \to \mathcal{Y}$ et $q : \mathcal{Z} \to \mathcal{Y}$
des $1$-morphismes de catégories fibrées en groupoïdes sur $(\Sch/S)_{fppf}$.
Si $p : \mathcal{X} \to \mathcal{Y}$ est formellement lisse sur les objets, il en
est de même du changement de base
$p' : \mathcal{X} \times_\mathcal{Y} \mathcal{Z} \to \mathcal{Z}$
de $p$ par $q$.
\end{lemma}

\begin{proof}
L'assertion est formelle. Soit $U \subset U'$ un épaississement du premier ordre
de schémas affines sur $S$, soit $z'$ un objet de $\mathcal{Z}$
au-dessus de $U'$, soit $w$ un objet de
$\mathcal{X} \times_\mathcal{Y} \mathcal{Z}$ au-dessus de $U$, et soit
$\delta : p'(w) \to z'|_U$ un isomorphisme.
On peut écrire
$w = (U, x, z, \alpha)$, où $x$ est un objet de $\mathcal{X}$ au-dessus de $U$,
$z$ un objet de $\mathcal{Z}$ au-dessus de $U$ et
$\alpha : p(x) \to q(z)$ un isomorphisme. Puisque $p'(w) = z$, on a
$\delta : z \to z'|_U$. Posons $y' = q(z')$ et
$\gamma = q(\delta) \circ \alpha : p(x) \to y'|_U$.
Comme $p$ est formellement lisse sur les objets, il existe un
objet $x'$ de $\mathcal{X}$ au-dessus de $U'$, ainsi que des
isomorphismes $\beta : x'|_U \to x$ et $\gamma' : p(x') \to y'$ tels que
(\ref{equation-formally-smooth}) soit commutatif. Considérons alors l'objet
$w' = (U', x', z', \gamma')$ de $\mathcal{X} \times_\mathcal{Y} \mathcal{Z}$
au-dessus de $U'$ et définissons les isomorphismes
$$
w'|_U = (U, x'|_U, z'|_U, \gamma'|_U)
\xrightarrow{(\beta, \delta^{-1})}
(U, x, z, \alpha) = w
$$
et
$$
p'(w') = z' \xrightarrow{\text{id}} z'.
$$
Ceux-ci se combinent en une solution du problème.
\end{proof}

\begin{lemma}
\label{lemma-composition-formally-smooth}
Soient $p : \mathcal{X} \to \mathcal{Y}$ et $q : \mathcal{Y} \to \mathcal{Z}$
des $1$-morphismes de catégories fibrées en groupoïdes sur $(\Sch/S)_{fppf}$.
Si $p$ et $q$ sont formellement lisses sur les objets, il en est de même de la composée
$q \circ p$.
\end{lemma}

\begin{proof}
L'assertion est formelle. Soit $U \subset U'$ un épaississement du premier ordre
de schémas affines sur $S$, soit $z'$ un objet de $\mathcal{Z}$
au-dessus de $U'$, soit $x$ un objet de $\mathcal{X}$ au-dessus de $U$,
et soit $\gamma : q(p(x)) \to z'|_U$ un isomorphisme. Comme $q$ est
formellement lisse sur les objets, il existe un objet
$y'$ de $\mathcal{Y}$ au-dessus de $U'$, un isomorphisme
$\beta : y'|_U \to p(x)$ et un isomorphisme $\gamma' : q(y') \to z'$
tels que (\ref{equation-formally-smooth}) soit commutatif. Comme $p$ est
formellement lisse sur les objets, il existe un objet
$x'$ de $\mathcal{X}$ au-dessus de $U'$, un isomorphisme
$\beta' : x'|_U \to x$ et un isomorphisme $\gamma'' : p(x') \to y'$
tels que (\ref{equation-formally-smooth}) soit commutatif.
La solution consiste à prendre $x'$ au-dessus de $U'$ avec l'isomorphisme
$$
q(p(x')) \xrightarrow{q(\gamma'')} q(y') \xrightarrow{\gamma'} z'
$$
et l'isomorphisme $\beta' : x'|_U \to x$. Nous omettons de vérifier
que (\ref{equation-formally-smooth}) est commutatif.
\end{proof}

\noindent
Notons que la classe des morphismes formellement lisses d'espaces algébriques est
stable par tout changement de base et locale sur le but pour la
topologie fpqc ; voir
More on Morphisms of Spaces,
lemmes \ref{spaces-more-morphisms-lemma-base-change-formally-smooth} et
\ref{spaces-more-morphisms-lemma-descending-property-formally-smooth}.
La condition (2) du lemme ci-dessous a donc un sens.

\begin{lemma}
\label{lemma-representable-by-spaces-formally-smooth}
Soit $p : \mathcal{X} \to \mathcal{Y}$ un $1$-morphisme de catégories
fibrées en groupoïdes sur $(\Sch/S)_{fppf}$. Si $p$ est
représentable par des espaces algébriques, les assertions suivantes sont équivalentes :
\begin{enumerate}
\item $p$ est formellement lisse sur les objets, et
\item $p$ est formellement lisse (voir
Algebraic Stacks,
définition \ref{algebraic-definition-relative-representable-property}).
\end{enumerate}
\end{lemma}

\begin{proof}
Supposons (2). Soit $U \subset U'$ un épaississement du premier ordre
de schémas affines sur $S$, soit $y'$ un objet de $\mathcal{Y}$
au-dessus de $U'$, soit $x$ un objet de $\mathcal{X}$ au-dessus de $U$,
et soit $\gamma : p(x) \to y'|_U$ un isomorphisme. Notons
$X_{y'}$ un espace algébrique sur $U'$ représentant le $2$-produit
fibré
$$
(\Sch/U')_{fppf} \times_{y', \mathcal{Y}, p} \mathcal{X}.
$$
Notons que $\xi = (U, U \to U', x, \gamma^{-1})$ définit un objet de
ce $2$-produit fibré au-dessus de $U$. Par le lemme de $2$-Yoneda, $\xi$ correspond
à un morphisme $f_\xi : U \to X_{y'}$ sur $U'$. Comme $X_{y'} \to U'$ est
formellement lisse par hypothèse, il existe un morphisme
$f' : U' \to X_{y'}$ tel que $f_\xi$ soit la composée de $f'$
et du morphisme $U \to U'$. En outre, le lemme de $2$-Yoneda montre que
$f'$ correspond à un objet $\xi' = (U', U' \to U', x', \alpha)$ du
$2$-produit fibré affiché au-dessus de $U'$, dont la restriction à
$U$ redonne $\xi$. En particulier, on obtient un isomorphisme
$\beta : x'|U \to x$. Notons que $\alpha : y' \to p(x')$.
Ainsi, prendre $x'$, l'isomorphisme
$\beta : x'|U \to x$ et l'isomorphisme
$\gamma' = \alpha^{-1} : p(x') \to y'$
fournit une solution du problème.

\medskip\noindent
Supposons (1). Choisissons un schéma $T$ et un $1$-morphisme
$y : (\Sch/T)_{fppf} \to \mathcal{Y}$. Soit
$X_y$ un espace algébrique sur $T$ représentant le $2$-produit fibré
$(\Sch/T)_{fppf} \times_{y, \mathcal{Y}, p} \mathcal{X}$.
Il faut montrer que $X_y \to T$ est formellement lisse.
Il suffit donc de montrer que, pour tout épaississement du premier ordre
$U \subset U'$ de schémas affines sur $T$,
$X_y(U') \to X_y(U)$ est surjective (il s'agit de morphismes dans la
catégorie des espaces algébriques sur $T$). Posons $y' = y|_{U'}$.
D'après le lemme de $2$-Yoneda, les morphismes $U \to X_y$ sur $T$ correspondent bijectivement
aux classes d'isomorphisme de couples $(x, \alpha)$, où $x$ est un objet
de $\mathcal{X}$ au-dessus de $U$ et $\alpha : y|_U \to p(x)$ un isomorphisme.
Bien entendu, donner $\alpha$ revient, à inversion près, à donner
un isomorphisme $\gamma : p(x) \to y'|_U$.
Il en va de même pour les morphismes $U' \to X_y$ sur $T$. Ainsi, (1) garantit
la surjectivité de $X_y(U') \to X_y(U)$
dans cette situation, ce qui conclut.
\end{proof}







\section{Surjectif sur les objets}
\label{section-formally-surjective}

\noindent
Soit $S$ un schéma. Soit $p : \mathcal{X} \to \mathcal{Y}$ un $1$-morphisme
de catégories fibrées en groupoïdes sur $(\Sch/S)_{fppf}$. Nous dirons que
$p$ est {\it surjectif sur les objets} si la condition suivante est satisfaite :
pour toute donnée formée de
\begin{enumerate}
\item un corps $k$ sur $S$, et
\item un objet $y$ de $\mathcal{Y}$ au-dessus de $\Spec(k)$,
\end{enumerate}
il existe une extension de corps $K/k$ sur $S$ et un
objet $x$ de $\mathcal{X}$ au-dessus de $\Spec(K)$
tel que $p(x) \cong y|_{\Spec(K)}$.

\begin{lemma}
\label{lemma-base-change-surjective}
Soient $p : \mathcal{X} \to \mathcal{Y}$ et $q : \mathcal{Z} \to \mathcal{Y}$
des $1$-morphismes de catégories fibrées en groupoïdes sur $(\Sch/S)_{fppf}$.
Si $p : \mathcal{X} \to \mathcal{Y}$ est surjectif sur les objets, il en
est de même du changement de base
$p' : \mathcal{X} \times_\mathcal{Y} \mathcal{Z} \to \mathcal{Z}$
de $p$ par $q$.
\end{lemma}

\begin{proof}
L'assertion est formelle. Soit $z$ un objet de $\mathcal{Z}$ au-dessus d'un corps $k$.
Comme $p$ est surjectif sur les objets, il existe une extension $K/k$,
un objet $x$ de $\mathcal{X}$ au-dessus de $K$ et un isomorphisme
$\alpha : p(x) \to q(z)|_{\Spec(K)}$. Alors
$w = (\Spec(K), x, z|_{\Spec(K)}, \alpha)$ est un objet de
$\mathcal{X} \times_\mathcal{Y} \mathcal{Z}$ au-dessus de $K$ tel que
$p'(w) = z|_{\Spec(K)}$.
\end{proof}

\begin{lemma}
\label{lemma-composition-surjective}
Soient $p : \mathcal{X} \to \mathcal{Y}$ et $q : \mathcal{Y} \to \mathcal{Z}$
des $1$-morphismes de catégories fibrées en groupoïdes sur $(\Sch/S)_{fppf}$.
Si $p$ et $q$ sont surjectifs sur les objets, il en est de même de la composée
$q \circ p$.
\end{lemma}

\begin{proof}
L'assertion est formelle. Soit $z$ un objet de $\mathcal{Z}$ au-dessus d'un corps $k$.
Comme $q$ est surjectif sur les objets, il existe une extension de corps $K/k$
et un objet $y$ de $\mathcal{Y}$ au-dessus de $K$ tels que
$q(y) \cong z|_{\Spec(K)}$. Comme $p$ est surjectif sur les objets, il
existe une extension de corps $L/K$ et un objet $x$ de $\mathcal{X}$
au-dessus de $L$ tel que $p(x) \cong y|_{\Spec(L)}$. Alors l'extension de corps
$L/k$ et l'objet $x$ de $\mathcal{X}$ au-dessus de $L$ vérifient
$q(p(x)) \cong z|_{\Spec(L)}$, comme voulu.
\end{proof}

\begin{lemma}
\label{lemma-representable-by-spaces-surjective}
Soit $p : \mathcal{X} \to \mathcal{Y}$ un $1$-morphisme de catégories
fibrées en groupoïdes sur $(\Sch/S)_{fppf}$. Si $p$ est
représentable par des espaces algébriques, les assertions suivantes sont équivalentes :
\begin{enumerate}
\item $p$ est surjectif sur les objets, et
\item $p$ est surjectif (voir
Algebraic Stacks,
définition \ref{algebraic-definition-relative-representable-property}).
\end{enumerate}
\end{lemma}

\begin{proof}
Supposons (2). Soit $k$ un corps et soit $y$ un objet de
$\mathcal{Y}$ au-dessus de $k$. Notons $X_y$ un espace algébrique sur $k$
représentant le $2$-produit fibré
$$
(\Sch/\Spec(k))_{fppf} \times_{y, \mathcal{Y}, p} \mathcal{X}.
$$
Puisque $p$ est supposé surjectif, $X_y$ n'est pas vide.
Il existe donc une extension de corps $K/k$ et un point $x$ de $X_y$
à valeurs dans $K$. Par le lemme de $2$-Yoneda, cela correspond à un objet
$x$ de $\mathcal{X}$ au-dessus de $K$, muni d'un isomorphisme
$p(x) \cong y|_{\Spec(K)}$ ; ainsi, (1) est satisfaite.

\medskip\noindent
Supposons (1). Choisissons un schéma $T$ et un $1$-morphisme
$y : (\Sch/T)_{fppf} \to \mathcal{Y}$. Soit
$X_y$ un espace algébrique sur $T$ représentant le $2$-produit fibré
$(\Sch/T)_{fppf} \times_{y, \mathcal{Y}, p} \mathcal{X}$.
Il faut montrer que $X_y \to T$ est surjectif. D'après
Morphisms of Spaces, définition \ref{spaces-morphisms-definition-surjective},
il faut montrer que $|X_y| \to |T|$ est surjectif.
Cela signifie exactement que, pour tout corps $k$ sur $T$ et tout
morphisme $t : \Spec(k) \to T$, il existe une extension de corps
$K/k$ et un morphisme $x : \Spec(K) \to X_y$ tels que
$$
\xymatrix{
\Spec(K) \ar[d] \ar[r]_x & X_y \ar[d] \\
\Spec(k) \ar[r]^t & T
}
$$
soit commutatif. D'après le lemme de $2$-Yoneda, cela signifie exactement qu'il faut trouver
$k \subset K$ et un objet $x$ de $\mathcal{X}$ au-dessus de $K$ tels que
$p(x) \cong t^*y|_{\Spec(K)}$. Ainsi, (1) garantit bien cette
existence, ce qui conclut.
\end{proof}












\section{Morphismes algébriques}
\label{section-algebraic}

\noindent
La notion suivante est parfois utile.

\begin{definition}
\label{definition-algebraic}
Soit $S$ un schéma. Soit $F : \mathcal{X} \to \mathcal{Y}$ un
$1$-morphisme de champs en groupoïdes sur $(\Sch/S)_{fppf}$.
On dit que $F$ est {\it algébrique} si, pour tout schéma $T$ et tout
objet $\xi$ de $\mathcal{Y}$ au-dessus de $T$, le $2$-produit fibré
$$
(\Sch/T)_{fppf} \times_{\xi, \mathcal{Y}} \mathcal{X}
$$
est un champ algébrique sur $S$.
\end{definition}

\noindent
Avec cette terminologie, on obtient le résultat suivant, qui généralise
Algebraic Stacks, lemme
\ref{algebraic-lemma-representable-morphism-to-algebraic}.

\begin{lemma}
\label{lemma-algebraic-morphism-to-algebraic}
Soit $S$ un schéma.
Soit $F : \mathcal{X} \to \mathcal{Y}$ un $1$-morphisme de
champs en groupoïdes sur $(\Sch/S)_{fppf}$. Si
\begin{enumerate}
\item $\mathcal{Y}$ est un champ algébrique, et
\item $F$ est algébrique (au sens ci-dessus),
\end{enumerate}
alors $\mathcal{X}$ est un champ algébrique.
\end{lemma}

\begin{proof}
D'après l'hypothèse (1), il existe un schéma $T$ et un objet
$\xi$ de $\mathcal{Y}$ au-dessus de $T$ tels que le
$1$-morphisme correspondant $\xi : (\Sch/T)_{fppf} \to \mathcal{Y}$
soit lisse et surjectif. Alors
$\mathcal{U} = (\Sch/T)_{fppf} \times_{\xi, \mathcal{Y}} \mathcal{X}$
est un champ algébrique par l'hypothèse (2).
Choisissons un schéma $U$ et un $1$-morphisme lisse surjectif
$(\Sch/U)_{fppf} \to \mathcal{U}$.
La projection $\mathcal{U} \longrightarrow \mathcal{X}$,
qui est le changement de base du morphisme
$\xi : (\Sch/T)_{fppf} \to \mathcal{Y}$,
est surjective et lisse ; voir
Algebraic Stacks, lemme
\ref{algebraic-lemma-base-change-representable-transformations-property}.
La composée
$(\Sch/U)_{fppf} \to \mathcal{U} \to \mathcal{X}$
est alors surjective et lisse comme composée de
morphismes surjectifs et lisses ; voir
Algebraic Stacks, lemme
\ref{algebraic-lemma-composition-representable-transformations-property}.
Ainsi, $\mathcal{X}$ est un champ algébrique d'après
Algebraic Stacks, lemme
\ref{algebraic-lemma-smooth-surjective-morphism-implies-algebraic}.
\end{proof}

\begin{lemma}
\label{lemma-map-from-algebraic}
Soit $S$ un schéma. Soit $F : \mathcal{X} \to \mathcal{Y}$ un $1$-morphisme
de champs en groupoïdes sur $(\Sch/S)_{fppf}$. Si $\mathcal{X}$ est un
champ algébrique et si $\Delta : \mathcal{Y} \to \mathcal{Y} \times \mathcal{Y}$
est représentable par des espaces algébriques, alors $F$ est algébrique.
\end{lemma}

\begin{proof}
Choisissons un champ en groupoïdes représentable $\mathcal{U}$ et un
$1$-morphisme lisse surjectif $\mathcal{U} \to \mathcal{X}$. Soit $T$ un schéma et
soit $\xi$ un objet de $\mathcal{Y}$ au-dessus de $T$. Le morphisme de
$2$-produits fibrés
$$
(\Sch/T)_{fppf} \times_{\xi, \mathcal{Y}} \mathcal{U}
\longrightarrow
(\Sch/T)_{fppf} \times_{\xi, \mathcal{Y}} \mathcal{X}
$$
est représentable par des espaces algébriques, surjectif et lisse comme
changement de base de $\mathcal{U} \to \mathcal{X}$ ; voir
Algebraic Stacks,
lemmes \ref{algebraic-lemma-base-change-representable-by-spaces} et
\ref{algebraic-lemma-base-change-representable-transformations-property}.
La condition imposée à la diagonale de $\mathcal{Y}$ montre que
la source de ce morphisme est représentable par un espace algébrique ; voir
Algebraic Stacks, lemme \ref{algebraic-lemma-representable-diagonal}.
La cible est donc un champ algébrique d'après
Algebraic Stacks,
lemme \ref{algebraic-lemma-smooth-surjective-morphism-implies-algebraic}.
\end{proof}

\begin{lemma}
\label{lemma-diagonals-and-algebraic-morphisms}
Soit $S$ un schéma. Soit $F : \mathcal{X} \to \mathcal{Y}$ un $1$-morphisme
de champs en groupoïdes sur $(\Sch/S)_{fppf}$.
Si $F$ est algébrique et si
$\Delta : \mathcal{Y} \to \mathcal{Y} \times \mathcal{Y}$
est représentable par des espaces algébriques, alors
$\Delta : \mathcal{X} \to \mathcal{X} \times \mathcal{X}$
est représentable par des espaces algébriques.
\end{lemma}

\begin{proof}
Supposons $F$ algébrique et
$\Delta : \mathcal{Y} \to \mathcal{Y} \times \mathcal{Y}$
représentable par des espaces algébriques.
Prenons un schéma $U$ sur $S$ et deux objets $x_1, x_2$ de
$\mathcal{X}$ au-dessus de $U$.
Il faut montrer que $\mathit{Isom}(x_1, x_2)$ est un espace algébrique
sur $U$ ; voir
Algebraic Stacks, lemme \ref{algebraic-lemma-representable-diagonal}.
Posons $y_i = F(x_i)$. On dispose d'un morphisme de faisceaux d'ensembles
$$
f : \mathit{Isom}(x_1, x_2) \to \mathit{Isom}(y_1, y_2)
$$
dont la cible est un espace algébrique par hypothèse.
Il suffit donc de montrer que $f$ est représentable par des
espaces algébriques ; voir Bootstrap, lemme
\ref{bootstrap-lemma-representable-by-spaces-over-space}.
On peut donc choisir un schéma $V$ sur $U$ et un
isomorphisme $\beta : y_{1, V} \to y_{2, V}$ ;
il faut montrer que le foncteur
$$
(\Sch/V)_{fppf} \to \textit{Sets},\quad
T/V \mapsto \{\alpha : x_{1, T} \to x_{2, T}
\text{ dans }\mathcal{X}_T \mid F(\alpha) = \beta|_T\}
$$
est un espace algébrique. Considérons les objets
$z_1 = (V, x_{1, V}, \text{id})$ et
$z_2 = (V, x_{2, V}, \beta)$ de
$$
\mathcal{Z} = (\Sch/V)_{fppf} \times_{y_{1, V}, \mathcal{Y}} \mathcal{X}
$$
On vérifie immédiatement que
le foncteur ci-dessus est égal à $\mathit{Isom}(z_1, z_2)$
sur $(\Sch/V)_{fppf}$. C'est donc un espace algébrique
par l'hypothèse que $F$ est algébrique (et par la définition
des champs algébriques).
\end{proof}
















\section{Espaces de sections}
\label{section-spaces-sections}

\noindent
Étant donnés des morphismes $W \to Z \to U$, on peut considérer le foncteur qui associe
à un schéma $U'$ sur $U$ l'ensemble des sections $\sigma : Z_{U'} \to W_{U'}$
du changement de base $W_{U'} \to Z_{U'}$ du morphisme $W \to Z$.
Dans cette section, nous démontrons quelques lemmes préliminaires sur ce foncteur.

\begin{lemma}
\label{lemma-surjection-space-of-sections}
Soit $Z \to U$ un morphisme fini de schémas.
Soit $W$ un espace algébrique et soit $W \to Z$ un
morphisme étale surjectif. Il existe alors un morphisme
étale surjectif $U' \to U$ et une section
$$
\sigma : Z_{U'} \to W_{U'}
$$
du morphisme $W_{U'} \to Z_{U'}$.
\end{lemma}

\begin{proof}
On peut choisir un schéma séparé $W'$ et un morphisme étale surjectif
$W' \to W$. Après avoir remplacé $W$ par $W'$, on peut donc supposer que $W$
est un schéma séparé. Écrivons $f : W \to Z$ et $\pi : Z \to U$.
Notons que $\pi \circ f : W \to U$ est séparé puisque
$W$ est séparé (voir
Schemes, lemme \ref{schemes-lemma-compose-after-separated}).
Soit $u \in U$ un point. Il suffit évidemment
de trouver un voisinage étale $(U', u')$ de $(U, u)$ tel qu'une
section $\sigma$ existe sur $U'$. Soient $z_1, \ldots, z_r$
les points de $Z$ situés au-dessus de $u$. Pour chaque $i$, choisissons un point
$w_i \in W$ d'image $z_i$. On peut choisir un voisinage étale
$(U', u') \to (U, u)$ tel que les conclusions de
More on Morphisms, Lemma
\ref{more-morphisms-lemma-etale-splits-off-quasi-finite-part-technical-variant}
soient satisfaites à la fois pour $Z \to U$ et les points $z_1, \ldots, z_r$,
et pour $W \to U$ et les points $w_1, \ldots, w_r$. Après avoir
remplacé $(U, u)$ par $(U', u')$ et renommé les objets, on peut donc supposer que
toutes les extensions de corps $\kappa(z_i)/\kappa(u)$ et
$\kappa(w_i)/\kappa(u)$ sont radicielles et, de plus,
qu'il existe des décompositions en sommes disjointes
$$
Z = V_1 \amalg \ldots \amalg V_r \amalg A, \quad
W = W_1 \amalg \ldots \amalg W_r \amalg B
$$
en sous-schémas ouverts et fermés,
avec $z_i \in V_i$, $w_i \in W_i$, et $V_i \to U$, $W_i \to U$ finis.
Après avoir remplacé $U$ par $U \setminus \pi(A)$, on peut supposer que
$A = \emptyset$, c'est-à-dire $Z = V_1 \amalg \ldots \amalg V_r$.
Après avoir remplacé $W_i$ par $W_i \cap f^{-1}(V_i)$ et
$B$ par $B \cup \bigcup W_i \cap f^{-1}(Z \setminus V_i)$,
on peut supposer que $f$ envoie $W_i$ dans $V_i$.
Alors $f_i = f|_{W_i} : W_i \to V_i$ est un morphisme entre schémas finis sur $U$,
donc il est fini (voir
Morphisms, lemme \ref{morphisms-lemma-finite-permanence}).
Il est aussi étale (par hypothèse),
$f_i^{-1}(\{z_i\}) = \{w_i\}$, et induit un isomorphisme des corps
résiduels $\kappa(z_i) = \kappa(w_i)$ (car les deux sont des extensions radicielles
de $\kappa(u)$, tandis que $\kappa(w_i)/\kappa(z_i)$
est séparable puisque $f$ est étale). D'après
\'Etale Morphisms, lemme \ref{etale-lemma-finite-etale-one-point}
le morphisme $f_i$ est donc un isomorphisme dans un voisinage $V_i'$ de
$z_i$. Comme $\pi : Z \to U$ est fermé, quitte à rétrécir $U$, on peut supposer
que $W_i \to V_i$ est un isomorphisme. Cela démontre le lemme.
\end{proof}

\begin{lemma}
\label{lemma-space-of-sections}
Soit $Z \to U$ un morphisme fini localement libre de schémas.
Soit $W$ un espace algébrique et soit $W \to Z$ un morphisme étale.
Alors le foncteur
$$
F : (\Sch/U)_{fppf}^{opp} \longrightarrow \textit{Sets},
$$
défini par la règle
$$
U' \longmapsto
F(U') =
\{\sigma : Z_{U'} \to W_{U'}\text{ section de }W_{U'} \to Z_{U'}\}
$$
est un espace algébrique et le morphisme $F \to U$ est étale.
\end{lemma}

\begin{proof}
Supposons d'abord que $W \to Z$ soit aussi séparé.
Soit $U'$ un schéma sur $U$ et soit $\sigma \in F(U')$. D'après
Morphisms of Spaces, lemme \ref{spaces-morphisms-lemma-section-immersion},
le morphisme $\sigma$ est une immersion fermée.
De plus, $\sigma$ est étale d'après
Properties of Spaces, lemme \ref{spaces-properties-lemma-etale-permanence}.
Ainsi, $\sigma$ est aussi une immersion ouverte ; voir
Morphisms of Spaces,
lemme \ref{spaces-morphisms-lemma-etale-universally-injective-open}.
Autrement dit, $Z_\sigma = \sigma(Z_{U'}) \subset W_{U'}$ est
un sous-espace ouvert tel que le morphisme $Z_\sigma \to Z_{U'}$
soit un isomorphisme. En particulier, le morphisme $Z_\sigma \to U'$
est fini. On obtient donc une transformation de foncteurs
$$
F \longrightarrow (W/U)_{fin}, \quad
\sigma \longmapsto (U' \to U, Z_\sigma)
$$
où $(W/U)_{fin}$ est la partie finie du morphisme $W \to U$
introduite dans
More on Groupoids in Spaces, section
\ref{spaces-more-groupoids-section-finite}.
Il est clair que cette transformation de foncteurs est injective (puisqu'on peut
retrouver $\sigma$ à partir de $Z_\sigma$ comme inverse de l'isomorphisme
$Z_\sigma \to Z_{U'}$). D'après
More on Groupoids in Spaces, proposition
\ref{spaces-more-groupoids-proposition-finite-algebraic-space}
nous savons que $(W/U)_{fin}$ est un espace algébrique étale sur $U$.
Pour achever la preuve dans ce cas, il suffit donc de montrer que
$F \to (W/U)_{fin}$ est représentable et est une immersion ouverte.
Pour cela, supposons donnés un morphisme de schémas $U' \to U$
et un sous-espace ouvert $Z' \subset W_{U'}$ tels que $Z' \to U'$
soit fini. Il suffit alors de montrer qu'il existe un
sous-schéma ouvert $U'' \subset U'$ tel qu'un morphisme
$T \to U'$ se factorise par $U''$ si et seulement si $Z' \times_{U'} T$
s'envoie isomorphiquement sur $Z \times_{U'} T$. Cela résulte de
More on Morphisms of Spaces, lemme
\ref{spaces-more-morphisms-lemma-where-isomorphism}
(on utilise ici que $Z \to U$ est plat et localement de présentation finie,
ainsi que fini).
Le lemme est donc démontré lorsque $W \to Z$ est séparé
et étale.

\medskip\noindent
Dans le cas général, choisissons un schéma séparé $W'$ et un morphisme
étale surjectif $W' \to W$. Notons que les morphismes $W' \to W$ et
$W' \to Z$ sont séparés puisque leur source est séparée. Notons $F'$ le
foncteur associé à $W' \to Z \to U$ comme dans le lemme. Dans le premier
paragraphe de la preuve, nous avons montré que $F'$ est représentable par un
espace algébrique étale sur $U$. D'après le
lemme \ref{lemma-surjection-space-of-sections},
l'application de foncteurs $F' \to F$ est surjective pour la topologie étale
sur $\Sch/U$. De plus, si $U'$ et $\sigma : Z_{U'} \to W_{U'}$
définissent un élément $\xi \in F(U')$, alors le produit fibré
$$
F'' = F' \times_{F, \xi} U'
$$
est le foncteur sur $\Sch/U'$ associé aux morphismes
$$
W'_{U'} \times_{W_{U'}, \sigma} Z_{U'} \to Z_{U'} \to U'.
$$
Comme le premier morphisme est séparé, étant le changement de base d'un morphisme
séparé, le résultat du premier paragraphe montre que $F''$ est un espace algébrique
étale sur $U'$. Il s'ensuit que $F' \to F$ est une transformation de foncteurs
étale surjective, représentable par des espaces algébriques. Ainsi, $F$ est
un espace algébrique d'après Bootstrap, théorème
\ref{bootstrap-theorem-final-bootstrap}. Comme $F' \to F$ est un
morphisme étale surjectif d'espaces algébriques, $F \to U$ est étale
puisque $F' \to U$ l'est.
\end{proof}









\section{Morphismes relatifs}
\label{section-relative-morphisms}

\noindent
Nous poursuivons la discussion commencée dans
More on Morphisms, section \ref{more-morphisms-section-relative-morphisms}.

\medskip\noindent
Soit $S$ un schéma. Soient $Z \to B$ et $X \to B$ des morphismes
d'espaces algébriques sur $S$. Pour tout schéma $T$, on peut considérer les couples
$(a, b)$, où $a : T \to B$
est un morphisme et $b : T \times_{a, B} Z \to T \times_{a, B} X$
un morphisme sur $T$. On a le diagramme
\begin{equation}
\label{equation-hom}
\vcenter{
\xymatrix{
T \times_{a, B} Z \ar[rd] \ar[rr]_b & &
T \times_{a, B} X \ar[ld] & Z \ar[rd] & & X \ar[ld] \\
& T \ar[rrr]^a & & & B
}
}
\end{equation}
Bien entendu, on peut aussi voir $b$ comme un morphisme
$b : T \times_{a, B} Z \to X$ tel que
$$
\xymatrix{
T \times_{a, B} Z \ar[r] \ar[d] \ar@/^1pc/[rrr]_-b &
Z \ar[rd] & & X \ar[ld] \\
T \ar[rr]^a & & B
}
$$
soit commutatif. Dans cette situation, on peut définir un foncteur
\begin{equation}
\label{equation-hom-functor}
\mathit{Mor}_B(Z, X) : (\Sch/S)^{opp} \longrightarrow \textit{Sets},
\quad
T \longmapsto \{(a, b)\text{ comme ci-dessus}\}
\end{equation}
Parfois, on le considère comme un foncteur défini sur la catégorie
des schémas sur $B$ ; dans ce cas, on omet $a$ de la notation.

\begin{lemma}
\label{lemma-hom-functor-sheaf}
Soit $S$ un schéma. Soient $Z \to B$ et $X \to B$ des morphismes
d'espaces algébriques sur $S$. Alors
\begin{enumerate}
\item $\mathit{Mor}_B(Z, X)$ est un faisceau sur
$(\Sch/S)_{fppf}$.
\item Si $T$ est un espace algébrique sur $S$, il existe une
bijection canonique
$$
\Mor_{\Sh((\Sch/S)_{fppf})}(T, \mathit{Mor}_B(Z, X))
=
\{(a, b)\text{ comme dans }(\ref{equation-hom})\}
$$
\end{enumerate}
\end{lemma}

\begin{proof}
Soit $T$ un espace algébrique sur $S$. Soit $\{T_i \to T\}$ un
recouvrement fppf de $T$ (comme dans
Topologies on Spaces, section \ref{spaces-topologies-section-fppf}).
Supposons que $(a_i, b_i) \in \mathit{Mor}_B(Z, X)(T_i)$ et
que $(a_i, b_i)|_{T_i \times_T T_j} = (a_j, b_j)|_{T_i \times_T T_j}$
pour tous $i, j$. D'après
Descent on Spaces,
lemme \ref{spaces-descent-lemma-fpqc-universal-effective-epimorphisms},
il existe un unique morphisme $a : T \to B$ tel que $a_i$ soit la
composée de $T_i \to T$ et de $a$. Alors
$\{T_i \times_{a_i, B} Z \to T \times_{a, B} Z\}$ est aussi un recouvrement
fppf, et le même lemme implique qu'il existe un unique morphisme
$b : T \times_{a, B} Z \to T \times_{a, B} X$ tel que $b_i$ soit la
composée de $T_i \times_{a_i, B} Z \to T \times_{a, B} Z$ et de $b$. Ainsi,
$(a, b) \in \mathit{Mor}_B(Z, X)(T)$ se restreint à $(a_i, b_i)$
sur $T_i$ pour tout $i$.

\medskip\noindent
Notons que le résultat du paragraphe précédent implique en particulier (1).

\medskip\noindent
Soit $T$ un espace algébrique sur $S$. Pour démontrer (2), nous allons
construire des applications réciproques entre les ensembles affichés. Dans la
suite, le mot « couple » désigne un couple $(a, b)$ s'insérant
dans (\ref{equation-hom}).

\medskip\noindent
Soit $v : T \to \mathit{Mor}_B(Z, X)$ une transformation naturelle.
Choisissons un schéma $U$ et un morphisme étale surjectif $p : U \to T$.
Alors $v(p) \in \mathit{Mor}_B(Z, X)(U)$ correspond à un couple $(a_U, b_U)$
sur $U$. Soit $R = U \times_T U$, de projections $t, s : R \to U$.
Comme $v$ est une transformation de foncteurs, les images inverses de
$(a_U, b_U)$ par $s$ et $t$ coïncident. Puisque $\{U \to T\}$ est un
recouvrement fppf, le résultat du premier paragraphe permet donc de
déduire qu'il existe un unique couple $(a, b)$ sur $T$.

\medskip\noindent
Réciproquement, soit $(a, b)$ un couple sur $T$.
Soient $U \to T$, $R = U \times_T U$ et $t, s : R \to U$ comme
ci-dessus. La restriction $(a, b)|_U$ donne alors naissance à une
transformation de foncteurs $v : h_U \to \mathit{Mor}_B(Z, X)$ par le
lemme de Yoneda
(Categories, lemme \ref{categories-lemma-yoneda}).
Comme les deux images inverses $s^*(a, b)|_U$ et $t^*(a, b)|_U$
sont égales, $v$ coégalise les deux applications
$h_t, h_s : h_R \to h_U$. Puisque $T = U/R$ est le faisceau quotient fppf d'après
Spaces, lemme \ref{spaces-lemma-space-presentation},
et puisque $\mathit{Mor}_B(Z, X)$ est un faisceau fppf d'après (1), on conclut
que $v$ se factorise par une application $T \to \mathit{Mor}_B(Z, X)$.

\medskip\noindent
Nous omettons de vérifier que les deux constructions ci-dessus sont
réciproques.
\end{proof}

\begin{lemma}
\label{lemma-base-change-hom-functor}
Soit $S$ un schéma. Soient $Z \to B$, $X \to B$ et $B' \to B$
des morphismes d'espaces algébriques sur $S$. Posons $Z' = B' \times_B Z$
et $X' = B' \times_B X$. Alors
$$
\mathit{Mor}_{B'}(Z', X')
=
B' \times_B \mathit{Mor}_B(Z, X)
$$
dans $\Sh((\Sch/S)_{fppf})$.
\end{lemma}

\begin{proof}
L'égalité des foncteurs résulte immédiatement des définitions.
L'égalité des faisceaux s'en déduit, car les deux membres sont des
faisceaux d'après le
lemme \ref{lemma-hom-functor-sheaf},
et parce qu'un produit fibré de faisceaux est égal au
produit fibré correspondant des préfaisceaux (c'est-à-dire des foncteurs).
\end{proof}

\begin{lemma}
\label{lemma-etale-covering-hom-functor}
Soit $S$ un schéma. Soient $Z \to B$ et $X' \to X \to B$ des morphismes
d'espaces algébriques sur $S$. Supposons que
\begin{enumerate}
\item $X' \to X$ soit étale, et
\item $Z \to B$ soit fini localement libre.
\end{enumerate}
Alors $\mathit{Mor}_B(Z, X') \to \mathit{Mor}_B(Z, X)$ est représentable
par des espaces algébriques et étale. Si $X' \to X$ est aussi surjectif,
alors $\mathit{Mor}_B(Z, X') \to \mathit{Mor}_B(Z, X)$ est surjectif.
\end{lemma}

\begin{proof}
Soit $U$ un schéma et soit $\xi = (a, b)$ un élément de
$\mathit{Mor}_B(Z, X)(U)$. Il faut montrer que le foncteur
$$
h_U \times_{\xi, \mathit{Mor}_B(Z, X)} \mathit{Mor}_B(Z, X')
$$
est représentable par un espace algébrique étale sur $U$. Posons
$Z_U = U \times_{a, B} Z$ et $W = Z_U \times_{b, X} X'$.
Alors $W \to Z_U \to U$ est comme dans le
lemme \ref{lemma-space-of-sections},
et le faisceau $F$ qui y est défini s'identifie au produit fibré
affiché ci-dessus. D'où la première assertion du lemme.
La seconde assertion résulte de celle-ci et du
lemme \ref{lemma-surjection-space-of-sections},
qui garantit que $F \to U$ est surjectif dans la situation ci-dessus.
\end{proof}

\begin{proposition}
\label{proposition-hom-functor-algebraic-space}
Soit $S$ un schéma. Soient $Z \to B$ et $X \to B$ des morphismes
d'espaces algébriques sur $S$. Si $Z \to B$ est fini localement libre,
alors $\mathit{Mor}_B(Z, X)$ est un espace algébrique.
\end{proposition}

\begin{proof}
Choisissons un schéma $B' = \coprod B'_i$, somme disjointe de
schémas affines $B'_i$, et un morphisme étale surjectif $B' \to B$.
On peut aussi supposer que $B'_i \times_B Z$ est le spectre d'un anneau
qui est un $\Gamma(B'_i, \mathcal{O}_{B'_i})$-module libre de type fini.
D'après le
lemme \ref{lemma-base-change-hom-functor}
et
Spaces, lemme
\ref{spaces-lemma-base-change-representable-transformations-property}
le morphisme $\mathit{Mor}_{B'}(Z', X') \to \mathit{Mor}_B(Z, X)$
est étale surjectif. D'après
Bootstrap, théorème \ref{bootstrap-theorem-final-bootstrap},
il suffit donc de démontrer la proposition lorsque $B = B'$ est une somme disjointe de
schémas affines $B'_i$ telle que chaque $B'_i \times_B Z$ soit fini libre
sur $B'_i$. Il suffit alors de démontrer le résultat après restriction
à chaque $B'_i$. On peut donc supposer que $B$ est affine et que
$\Gamma(Z, \mathcal{O}_Z)$ est un $\Gamma(B, \mathcal{O}_B)$-module libre de type fini.

\medskip\noindent
Choisissons un schéma $X'$ qui soit une somme disjointe de schémas affines et
un morphisme étale surjectif $X' \to X$. D'après le
lemme \ref{lemma-etale-covering-hom-functor},
le morphisme $\mathit{Mor}_B(Z, X') \to \mathit{Mor}_B(Z, X)$
est représentable par des espaces algébriques, étale et surjectif.
D'après
Bootstrap, théorème \ref{bootstrap-theorem-final-bootstrap},
il suffit donc de démontrer la proposition lorsque $X$ est une somme disjointe
de schémas affines. On se ramène ainsi au cas traité dans le paragraphe
suivant.

\medskip\noindent
Supposons que $X = \coprod_{i \in I} X_i$ soit une somme disjointe de schémas
affines, que $B$ soit affine et que $\Gamma(Z, \mathcal{O}_Z)$ soit un
$\Gamma(B, \mathcal{O}_B)$-module libre de type fini. Pour toute partie finie
$E \subset I$, posons
$$
F_E = \mathit{Mor}_B(Z, \coprod\nolimits_{i \in E} X_i).
$$
D'après More on Morphisms,
lemme \ref{more-morphisms-lemma-hom-from-finite-free-into-affine},
$F_E$ est un espace algébrique. Considérons le morphisme
$$
\coprod\nolimits_{E \subset I\text{ fini}} F_E
\longrightarrow
\mathit{Mor}_B(Z, X)
$$
Chacun des morphismes
$F_E \to \mathit{Mor}_B(Z, X)$ est une immersion ouverte, car son image est
simplement le lieu paramétrant les couples $(a, b)$ pour lesquels $b$ se factorise par
le sous-schéma ouvert $\coprod\nolimits_{i \in E} X_i$ de $X$. De plus,
si $T$ est quasi-compact, alors, pour tout couple $(a, b)$, l'image
de $b$ est contenue dans $\coprod\nolimits_{i \in E} X_i$ pour une certaine partie
finie $E \subset I$. La flèche affichée est donc en fait un
recouvrement ouvert, ce qui conclut\footnote{Sous réserve de
quelques arguments ensemblistes. Il faut en effet montrer que
$\coprod F_E$ est un espace algébrique. Cela résulte des inégalités
$|I| \leq \text{size}(X)$ et $\text{size}(F_E) \leq \text{size}(X)$,
la seconde découlant de la description explicite de $F_E$ dans la preuve de
More on Morphisms,
lemme \ref{more-morphisms-lemma-hom-from-finite-free-into-affine}.
Nous omettons certains détails.}, d'après
Spaces, lemme \ref{spaces-lemma-glueing-algebraic-spaces}.
\end{proof}










\section{Restriction des scalaires}
\label{section-restriction-of-scalars}

\noindent
Supposons que $X \to Z \to B$ soient des morphismes d'espaces algébriques sur $S$.
Pour tout schéma $T$, on peut considérer les couples $(a, b)$, où $a : T \to B$
est un morphisme et $b : T \times_{a, B} Z \to X$ un morphisme sur $Z$.
On a le diagramme
\begin{equation}
\label{equation-pairs}
\vcenter{
\xymatrix{
& X \ar[d] \\
T \times_{a, B} Z \ar[d] \ar[ru]^b \ar[r] & Z \ar[d] \\
T \ar[r]^a & B
}
}
\end{equation}
Dans cette situation, on peut définir un
foncteur
\begin{equation}
\label{equation-restriction-of-scalars}
\text{Res}_{Z/B}(X) : (\Sch/S)^{opp} \longrightarrow \textit{Sets},
\quad
T \longmapsto \{(a, b)\text{ comme ci-dessus}\}
\end{equation}
Parfois, on le considère comme un foncteur défini sur la catégorie
des schémas sur $B$ ; dans ce cas, on omet $a$ de la notation.

\begin{lemma}
\label{lemma-restriction-of-scalars-sheaf}
Soit $S$ un schéma. Soient $X \to Z \to B$ des morphismes
d'espaces algébriques sur $S$. Alors
\begin{enumerate}
\item $\text{Res}_{Z/B}(X)$ est un faisceau sur
$(\Sch/S)_{fppf}$.
\item Si $T$ est un espace algébrique sur $S$, il existe une
bijection canonique
$$
\Mor_{\Sh((\Sch/S)_{fppf})}(T, \text{Res}_{Z/B}(X))
=
\{(a, b)\text{ comme dans }(\ref{equation-pairs})\}
$$
\end{enumerate}
\end{lemma}

\begin{proof}
Soit $T$ un espace algébrique sur $S$. Soit $\{T_i \to T\}$ un
recouvrement fppf de $T$ (comme dans
Topologies on Spaces, section \ref{spaces-topologies-section-fppf}).
Supposons que $(a_i, b_i) \in \text{Res}_{Z/B}(X)(T_i)$ et
que $(a_i, b_i)|_{T_i \times_T T_j} = (a_j, b_j)|_{T_i \times_T T_j}$
pour tous $i, j$. D'après
Descent on Spaces,
lemme \ref{spaces-descent-lemma-fpqc-universal-effective-epimorphisms},
il existe un unique morphisme $a : T \to B$ tel que $a_i$ soit la
composée de $T_i \to T$ et de $a$. Alors
$\{T_i \times_{a_i, B} Z \to T \times_{a, B} Z\}$ est aussi un recouvrement
fppf, et le même lemme implique qu'il existe un unique morphisme
$b : T \times_{a, B} Z \to X$ tel que $b_i$ soit la composée
de $T_i \times_{a_i, B} Z \to T \times_{a, B} Z$ et de $b$. Ainsi,
$(a, b) \in \text{Res}_{Z/B}(X)(T)$ se restreint à $(a_i, b_i)$
sur $T_i$ pour tout $i$.

\medskip\noindent
Notons que le résultat du paragraphe précédent implique en particulier (1).

\medskip\noindent
Soit $T$ un espace algébrique sur $S$. Pour démontrer (2), nous allons
construire des applications réciproques entre les ensembles affichés. Dans la
suite, le mot « couple » désigne un couple $(a, b)$ s'insérant
dans (\ref{equation-pairs}).

\medskip\noindent
Soit $v : T \to \text{Res}_{Z/B}(X)$ une transformation naturelle.
Choisissons un schéma $U$ et un morphisme étale surjectif $p : U \to T$.
Alors $v(p) \in \text{Res}_{Z/B}(X)(U)$ correspond à un couple $(a_U, b_U)$
sur $U$. Soit $R = U \times_T U$, de projections $t, s : R \to U$.
Comme $v$ est une transformation de foncteurs, les images inverses de
$(a_U, b_U)$ par $s$ et $t$ coïncident. Puisque $\{U \to T\}$ est un
recouvrement fppf, le résultat du premier paragraphe permet donc de
déduire qu'il existe un unique couple $(a, b)$ sur $T$.

\medskip\noindent
Réciproquement, soit $(a, b)$ un couple sur $T$.
Soient $U \to T$, $R = U \times_T U$ et $t, s : R \to U$ comme
ci-dessus. La restriction $(a, b)|_U$ donne alors naissance à une
transformation de foncteurs $v : h_U \to \text{Res}_{Z/B}(X)$ par le
lemme de Yoneda
(Categories, lemme \ref{categories-lemma-yoneda}).
Comme les deux images inverses $s^*(a, b)|_U$ et $t^*(a, b)|_U$
sont égales, $v$ coégalise les deux applications
$h_t, h_s : h_R \to h_U$. Puisque $T = U/R$ est le faisceau quotient fppf d'après
Spaces, lemme \ref{spaces-lemma-space-presentation},
et puisque $\text{Res}_{Z/B}(X)$ est un faisceau fppf d'après (1), on conclut
que $v$ se factorise par une application $T \to \text{Res}_{Z/B}(X)$.

\medskip\noindent
Nous omettons de vérifier que les deux constructions ci-dessus sont
réciproques.
\end{proof}

\noindent
Bien entendu, le faisceau $\text{Res}_{Z/B}(X)$ est muni d'une transformation
naturelle de foncteurs $\text{Res}_{Z/B}(X) \to B$. Nous l'utiliserons sans autre
mention dans la suite.

\begin{lemma}
\label{lemma-etale-base-change-restriction-of-scalars}
Soit $S$ un schéma. Soient $X \to Z \to B$ et $B' \to B$
des morphismes d'espaces algébriques sur $S$.
Posons $Z' = B' \times_B Z$ et $X' = B' \times_B X$. Alors
$$
\text{Res}_{Z'/B'}(X')
=
B' \times_B \text{Res}_{Z/B}(X)
$$
dans $\Sh((\Sch/S)_{fppf})$.
\end{lemma}

\begin{proof}
L'égalité des foncteurs résulte immédiatement des définitions.
L'égalité des faisceaux s'en déduit, car les deux membres sont des
faisceaux d'après le
lemme \ref{lemma-restriction-of-scalars-sheaf},
et parce qu'un produit fibré de faisceaux est égal au
produit fibré correspondant des préfaisceaux (c'est-à-dire des foncteurs).
\end{proof}

\begin{lemma}
\label{lemma-etale-covering-restriction-of-scalars}
Soit $S$ un schéma. Soient $X' \to X \to Z \to B$ des morphismes
d'espaces algébriques sur $S$. Supposons que
\begin{enumerate}
\item $X' \to X$ soit étale, et
\item $Z \to B$ soit fini localement libre.
\end{enumerate}
Alors $\text{Res}_{Z/B}(X') \to \text{Res}_{Z/B}(X)$ est représentable
par des espaces algébriques et étale. Si $X' \to X$ est aussi surjectif,
alors $\text{Res}_{Z/B}(X') \to \text{Res}_{Z/B}(X)$ est surjectif.
\end{lemma}

\begin{proof}
Soit $U$ un schéma et soit $\xi = (a, b)$ un élément de
$\text{Res}_{Z/B}(X)(U)$. Il faut montrer que le foncteur
$$
h_U \times_{\xi, \text{Res}_{Z/B}(X)} \text{Res}_{Z/B}(X')
$$
est représentable par un espace algébrique étale sur $U$. Posons
$Z_U = U \times_{a, B} Z$ et $W = Z_U \times_{b, X} X'$.
Alors $W \to Z_U \to U$ est comme dans le
lemme \ref{lemma-space-of-sections},
et le faisceau $F$ qui y est défini s'identifie au produit fibré
affiché ci-dessus. D'où la première assertion du lemme.
La seconde assertion résulte de celle-ci et du
lemme \ref{lemma-surjection-space-of-sections},
qui garantit que $F \to U$ est surjectif dans la situation ci-dessus.
\end{proof}

\noindent
On peut maintenant utiliser les lemmes ci-dessus pour montrer que $\text{Res}_{Z/B}(X)$
est un espace algébrique dès que $Z \to B$ est fini localement libre, d'une manière
presque identique à la preuve que $\mathit{Mor}_B(Z, X)$ est un
espace algébrique ; voir la
proposition \ref{proposition-hom-functor-algebraic-space}.
Nous déduirons plutôt directement ce résultat du lemme suivant
et du fait que $\mathit{Mor}_B(Z, X)$ est un espace algébrique.

\begin{lemma}
\label{lemma-fibre-diagram}
Soit $S$ un schéma. Soient $X \to Z \to B$ des morphismes
d'espaces algébriques sur $S$. Le diagramme suivant
$$
\xymatrix{
\mathit{Mor}_B(Z, X) \ar[r] & \mathit{Mor}_B(Z, Z) \\
\text{Res}_{Z/B}(X) \ar[r] \ar[u] & B \ar[u]_{\text{id}_Z}
}
$$
est un diagramme cartésien de faisceaux sur $(\Sch/S)_{fppf}$.
\end{lemma}

\begin{proof}
La preuve est omise. Indication : exercice sur le point de vue fonctoriel en
géométrie algébrique.
\end{proof}

\begin{proposition}
\label{proposition-restriction-of-scalars-algebraic-space}
Soit $S$ un schéma. Soient $X \to Z \to B$ des morphismes
d'espaces algébriques sur $S$. Si $Z \to B$ est fini localement libre,
alors $\text{Res}_{Z/B}(X)$ est un espace algébrique.
\end{proposition}

\begin{proof}
D'après la
proposition \ref{proposition-hom-functor-algebraic-space},
les foncteurs $\mathit{Mor}_B(Z, X)$ et $\mathit{Mor}_B(Z, Z)$
sont des espaces algébriques. L'assertion résulte donc du diagramme cartésien
du
lemme \ref{lemma-fibre-diagram}
et du fait que les produits fibrés d'espaces algébriques existent et
sont donnés par le produit fibré dans la catégorie sous-jacente des
faisceaux d'ensembles (voir
Spaces, lemme
\ref{spaces-lemma-fibre-product-spaces-over-sheaf-with-representable-diagonal}).
\end{proof}






\section{Champs de Hilbert finis}
\label{section-finite-hilbert-stacks}

\noindent
Dans cette section, nous démontrons quelques résultats sur les champs de
Hilbert finis $\mathcal{H}_d(\mathcal{X}/\mathcal{Y})$
introduits dans
Examples of Stacks, section \ref{examples-stacks-section-hilbert-d-stack}.

\begin{lemma}
\label{lemma-map-hilbert}
Considérons un diagramme $2$-commutatif
$$
\xymatrix{
\mathcal{X}' \ar[r]_G \ar[d]_{F'} & \mathcal{X} \ar[d]^F \\
\mathcal{Y}' \ar[r]^H & \mathcal{Y}
}
$$
de champs en groupoïdes sur $(\Sch/S)_{fppf}$, muni d'un
$2$-isomorphisme $\gamma : H \circ F' \to F \circ G$. Dans cette situation, on
obtient un $1$-morphisme canonique
$\mathcal{H}_d(\mathcal{X}'/\mathcal{Y}') \to
\mathcal{H}_d(\mathcal{X}/\mathcal{Y})$.
Ce morphisme est compatible avec les $1$-morphismes d'oubli de
Examples of Stacks,
équation (\ref{examples-stacks-equation-diagram-hilbert-d-stack}).
\end{lemma}

\begin{proof}
On envoie l'objet $(U, Z, y', x', \alpha')$ sur l'objet
$(U, Z, H(y'), G(x'), \gamma \star \text{id}_H \star \alpha')$,
où $\star$ désigne la composition horizontale des $2$-morphismes ; voir
Categories, définition \ref{categories-definition-horizontal-composition}.
Au morphisme
$(f, g, b, a) :
(U_1, Z_1, y_1', x_1', \alpha_1') \to (U_2, Z_2, y_2', x_2', \alpha_2')$
on associe
$(f, g, H(b), G(a))$.
Nous omettons de vérifier que cela définit un foncteur entre catégories sur
$(\Sch/S)_{fppf}$.
\end{proof}

\begin{lemma}
\label{lemma-cartesian-map-hilbert}
Dans la situation du
lemme \ref{lemma-map-hilbert},
supposons que le carré donné soit $2$-cartésien. Alors le diagramme
$$
\xymatrix{
\mathcal{H}_d(\mathcal{X}'/\mathcal{Y}') \ar[r] \ar[d] &
\mathcal{H}_d(\mathcal{X}/\mathcal{Y}) \ar[d] \\
\mathcal{Y}' \ar[r] &
\mathcal{Y}
}
$$
est $2$-cartésien.
\end{lemma}

\begin{proof}
Le lemme \ref{lemma-map-hilbert}
fournit un diagramme $2$-commutatif et, par conséquent,
un $1$-morphisme (c'est-à-dire un foncteur)
$$
\mathcal{H}_d(\mathcal{X}'/\mathcal{Y}')
\longrightarrow
\mathcal{Y}' \times_\mathcal{Y} \mathcal{H}_d(\mathcal{X}/\mathcal{Y})
$$
Indiquons pourquoi ce foncteur est essentiellement surjectif. Un objet
de la catégorie du membre de droite consiste en un schéma $U$ sur $S$,
un objet $y'$ de $\mathcal{Y}'_U$, un objet $(U, Z, y, x, \alpha)$
de $\mathcal{H}_d(\mathcal{X}/\mathcal{Y})$ au-dessus de $U$ et un isomorphisme
$H(y') \to y$ dans $\mathcal{Y}_U$. L'hypothèse signifie exactement qu'il
existe un objet $x'$ de $\mathcal{X}'_Z$ muni d'isomorphismes
$G(x') \cong x$ et $\alpha' : y'|_Z \to F'(x')$ compatibles
avec $\alpha$. Alors $(U, Z, y', x', \alpha')$ est un
objet de $\mathcal{H}_d(\mathcal{X}'/\mathcal{Y}')$ au-dessus de $U$.
Nous omettons les détails.
\end{proof}

\begin{lemma}
\label{lemma-etale-covering-hilbert}
Dans la situation du
lemme \ref{lemma-map-hilbert},
supposons que
\begin{enumerate}
\item $\mathcal{Y}' = \mathcal{Y}$ et $H = \text{id}_\mathcal{Y}$,
\item $G$ soit représentable par des espaces algébriques et étale.
\end{enumerate}
Alors $\mathcal{H}_d(\mathcal{X}'/\mathcal{Y}) \to
\mathcal{H}_d(\mathcal{X}/\mathcal{Y})$ est représentable par
des espaces algébriques et étale.
Si $G$ est en outre surjectif, alors
$\mathcal{H}_d(\mathcal{X}'/\mathcal{Y}) \to
\mathcal{H}_d(\mathcal{X}/\mathcal{Y})$ est surjectif.
\end{lemma}

\begin{proof}
Soit $U$ un schéma et soit $\xi = (U, Z, y, x, \alpha)$ un objet de
$\mathcal{H}_d(\mathcal{X}/\mathcal{Y})$ au-dessus de $U$.
Il faut montrer que le $2$-produit fibré
\begin{equation}
\label{equation-to-show}
(\Sch/U)_{fppf}
\times_{\xi, \mathcal{H}_d(\mathcal{X}/\mathcal{Y})}
\mathcal{H}_d(\mathcal{X}'/\mathcal{Y})
\end{equation}
est représentable par un espace algébrique étale sur $U$.
Un objet de celui-ci au-dessus de $U'$ correspond à un objet
$x'$ de la catégorie fibre de $\mathcal{X}'$ au-dessus de $Z_{U'}$
tel que $G(x') \cong x|_{Z_{U'}}$.
Par hypothèse, le $2$-produit fibré
$$
(\Sch/Z)_{fppf} \times_{x, \mathcal{X}} \mathcal{X}'
$$
est représentable par un espace algébrique $W$ tel que la projection
$W \to Z$ soit étale. Alors (\ref{equation-to-show})
est représenté par l'espace algébrique $F$ qui paramètre les sections de
$W \to Z$ sur $U$, introduit dans le
lemme \ref{lemma-space-of-sections}.
Comme $F \to U$ est étale, on conclut que
$\mathcal{H}_d(\mathcal{X}'/\mathcal{Y}) \to
\mathcal{H}_d(\mathcal{X}/\mathcal{Y})$ est représentable par
des espaces algébriques et étale.
Enfin, si $\mathcal{X}' \to \mathcal{X}$ est aussi surjectif,
alors $W \to Z$ est surjectif, et donc $F \to U$ l'est d'après le
lemme \ref{lemma-surjection-space-of-sections}.
Ainsi, dans ce cas,
$\mathcal{H}_d(\mathcal{X}'/\mathcal{Y}) \to
\mathcal{H}_d(\mathcal{X}/\mathcal{Y})$ est aussi surjectif.
\end{proof}

\begin{lemma}
\label{lemma-etale-map-hilbert}
Dans la situation du
lemme \ref{lemma-map-hilbert},
supposons que $G$ et $H$ soient représentables par des espaces algébriques et étales.
Alors $\mathcal{H}_d(\mathcal{X}'/\mathcal{Y}') \to
\mathcal{H}_d(\mathcal{X}/\mathcal{Y})$ est représentable par
des espaces algébriques et étale.
Si, de plus, $H$ est surjectif et le foncteur induit
$\mathcal{X}' \to \mathcal{Y}' \times_\mathcal{Y} \mathcal{X}$
est surjectif, alors
$\mathcal{H}_d(\mathcal{X}'/\mathcal{Y}') \to
\mathcal{H}_d(\mathcal{X}/\mathcal{Y})$ est surjectif.
\end{lemma}

\begin{proof}
Posons $\mathcal{X}'' = \mathcal{Y}' \times_\mathcal{Y} \mathcal{X}$. D'après le
lemme \ref{lemma-etale-permanence},
le $1$-morphisme $\mathcal{X}' \to \mathcal{X}''$ est représentable par
des espaces algébriques et étale (en particulier, la condition de la seconde
assertion du lemme selon laquelle $\mathcal{X}' \to \mathcal{X}''$ est surjectif
a un sens). On obtient un diagramme $2$-commutatif
$$
\xymatrix{
\mathcal{X}' \ar[r] \ar[d] &
\mathcal{X}'' \ar[r] \ar[d] &
\mathcal{X} \ar[d] \\
\mathcal{Y}' \ar[r] &
\mathcal{Y}' \ar[r] &
\mathcal{Y}
}
$$
Il résulte du
lemme \ref{lemma-cartesian-map-hilbert}
que $\mathcal{H}_d(\mathcal{X}''/\mathcal{Y}')$ est le changement de base
de $\mathcal{H}_d(\mathcal{X}/\mathcal{Y})$ par $\mathcal{Y}' \to \mathcal{Y}$.
En particulier,
$\mathcal{H}_d(\mathcal{X}''/\mathcal{Y}') \to
\mathcal{H}_d(\mathcal{X}/\mathcal{Y})$ est
représentable par des espaces algébriques et étale ; voir
Algebraic Stacks, lemme
\ref{algebraic-lemma-base-change-representable-transformations-property}.
De plus, il est aussi surjectif si $H$ l'est.
Par conséquent, il suffit de montrer que
le résultat vaut pour le carré de gauche du diagramme.
On se ramène ainsi au cas où $\mathcal{Y}' = \mathcal{Y}$,
qui fait l'objet du
lemme \ref{lemma-etale-covering-hilbert}.
\end{proof}

\begin{lemma}
\label{lemma-relative-hilbert}
Soit $F : \mathcal{X} \to \mathcal{Y}$ un $1$-morphisme de champs en groupoïdes
sur $(\Sch/S)_{fppf}$. Supposons que
$\Delta : \mathcal{Y} \to \mathcal{Y} \times \mathcal{Y}$
soit représentable par des espaces algébriques. Alors
$$
\mathcal{H}_d(\mathcal{X}/\mathcal{Y})
\longrightarrow
\mathcal{H}_d(\mathcal{X}) \times \mathcal{Y}
$$
qui figure dans
Examples of Stacks, équation
(\ref{examples-stacks-equation-diagram-hilbert-d-stack}),
est représentable par des espaces algébriques.
\end{lemma}

\begin{proof}
Soit $U$ un schéma et soit $\xi = (U, Z, p, x, 1)$ un objet de
$\mathcal{H}_d(\mathcal{X}) = \mathcal{H}_d(\mathcal{X}/S)$ au-dessus de $U$.
Ici, $p$ est simplement le morphisme structural de $U$.
La cinquième composante $1$ existe et est unique
puisque tout est au-dessus de $S$.
Soit aussi $y$ un objet de $\mathcal{Y}$ au-dessus de $U$.
Il faut montrer que le $2$-produit fibré
\begin{equation}
\label{equation-res-isom}
(\Sch/U)_{fppf}
\times_{\xi \times y, \mathcal{H}_d(\mathcal{X}) \times \mathcal{Y}}
\mathcal{H}_d(\mathcal{X}/\mathcal{Y})
\end{equation}
est représentable par un espace algébrique. Pour l'expliquer,
introduisons
$$
I = \mathit{Isom}_\mathcal{Y}(y|_Z, F(x))
$$
qui est un espace algébrique sur $Z$ par hypothèse. Soit $a : U' \to U$
un schéma sur $U$. Que signifie se donner un objet de la catégorie
fibre de (\ref{equation-res-isom}) au-dessus de $U'$ ? Cela signifie que l'on
dispose d'un objet $\xi' = (U', Z', y', x', \alpha')$ de
$\mathcal{H}_d(\mathcal{X}/\mathcal{Y})$ au-dessus de $U'$ et d'isomorphismes
$(U', Z', p', x', 1) \cong (U, Z, p, x, 1)|_{U'}$ et
$y' \cong y|_{U'}$. Ainsi, $\xi'$ est isomorphe à
$(U', U' \times_{a, U} Z, a^*y, x|_{U' \times_{a, U} Z}, \alpha)$
pour un certain morphisme
$$
\alpha :
a^*y|_{U' \times_{a, U} Z}
\longrightarrow
F(x|_{U' \times_{a, U} Z})
$$
dans la catégorie fibre de $\mathcal{Y}$ au-dessus de $U' \times_{a, U} Z$. Ainsi,
on peut voir $\alpha$ comme un morphisme $b : U' \times_{a, U} Z \to I$.
On voit de cette manière que (\ref{equation-res-isom})
est représenté par $\text{Res}_{Z/U}(I)$, qui est un espace algébrique d'après la
proposition \ref{proposition-restriction-of-scalars-algebraic-space}.
\end{proof}

\noindent
Le lemme suivant est une généralisation (partielle) du
lemme \ref{lemma-etale-covering-hilbert}.

\begin{lemma}
\label{lemma-representable-on-top}
Soient $F : \mathcal{X} \to \mathcal{Y}$ et $G : \mathcal{X}' \to \mathcal{X}$
des $1$-morphismes de champs en groupoïdes sur $(\Sch/S)_{fppf}$.
Si $G$ est représentable par des espaces algébriques, alors le $1$-morphisme
$$
\mathcal{H}_d(\mathcal{X}'/\mathcal{Y})
\longrightarrow
\mathcal{H}_d(\mathcal{X}/\mathcal{Y})
$$
est représentable par des espaces algébriques.
\end{lemma}

\begin{proof}
Soit $U$ un schéma et soit $\xi = (U, Z, y, x, \alpha)$ un objet de
$\mathcal{H}_d(\mathcal{X}/\mathcal{Y})$ au-dessus de $U$.
Il faut montrer que le $2$-produit fibré
\begin{equation}
\label{equation-to-show-again}
(\Sch/U)_{fppf}
\times_{\xi, \mathcal{H}_d(\mathcal{X}/\mathcal{Y})}
\mathcal{H}_d(\mathcal{X}'/\mathcal{Y})
\end{equation}
est représentable par un espace algébrique sur $U$.
Un objet de celui-ci au-dessus de $a : U' \to U$ correspond à un objet
$x'$ de $\mathcal{X}'$ au-dessus de $U' \times_{a, U} Z$ tel que
$G(x') \cong x|_{U' \times_{a, U} Z}$. Par hypothèse, le $2$-produit fibré
$$
(\Sch/Z)_{fppf} \times_{x, \mathcal{X}} \mathcal{X}'
$$
est représentable par un espace algébrique $W$ sur $Z$. Il s'ensuit que
(\ref{equation-to-show-again}) est représenté par $\text{Res}_{Z/U}(W)$,
qui est un espace algébrique d'après la
proposition \ref{proposition-restriction-of-scalars-algebraic-space}.
\end{proof}

\begin{lemma}
\label{lemma-limit-preserving}
Soit $F : \mathcal{X} \to \mathcal{Y}$ un $1$-morphisme de champs en groupoïdes
sur $(\Sch/S)_{fppf}$. Supposons que $F$ soit représentable par des espaces
algébriques et localement de présentation finie. Alors
$$
p : \mathcal{H}_d(\mathcal{X}/\mathcal{Y}) \to \mathcal{Y}
$$
commute aux limites sur les objets.
\end{lemma}

\begin{proof}
Cela signifie qu'il faut montrer ceci : étant donnés
\begin{enumerate}
\item un schéma affine $U = \lim_i U_i$ écrit comme la
limite projective filtrante de schémas affines $U_i$ sur $S$,
\item un objet $y_i$ de $\mathcal{Y}$ au-dessus de $U_i$ pour un certain $i$, et
\item un objet $\Xi = (U, Z, y, x, \alpha)$ de
$\mathcal{H}_d(\mathcal{X}/\mathcal{Y})$
au-dessus de $U$ tel que $y = y_i|_U$,
\end{enumerate}
il existe $i' \geq i$ et un objet
$\Xi_{i'} = (U_{i'}, Z_{i'}, y_{i'}, x_{i'}, \alpha_{i'})$ de
$\mathcal{H}_d(\mathcal{X}/\mathcal{Y})$ au-dessus de $U_{i'}$ tels que
$\Xi_{i'}|_U = \Xi$ et $y_{i'} = y_i|_{U_{i'}}$.
En effet, ces deux dernières égalités assurent la commutativité de
(\ref{equation-limit-preserving}).

\medskip\noindent
Soit $X_{y_i} \to U_i$ un espace algébrique représentant le $2$-produit
fibré
$$
(\Sch/U_i)_{fppf} \times_{y_i, \mathcal{Y}, F} \mathcal{X}.
$$
Notons que $X_{y_i} \to U_i$ est localement de présentation finie par notre
hypothèse sur $F$. Il est clair que
$\xi = (Z, Z \to U_i, x, \alpha)$ est un objet du $2$-produit fibré
affiché ci-dessus ; ainsi, $\xi$ donne naissance à un morphisme
$f_\xi : Z \to X_{y_i}$ d'espaces algébriques sur $U_i$
(puisque $X_{y_i}$ est le foncteur des classes d'isomorphisme d'objets de
$(\Sch/U_i)_{fppf} \times_{y_i, \mathcal{Y}, F} \mathcal{X}$ ; voir
Algebraic Stacks,
lemme \ref{algebraic-lemma-characterize-representable-by-space}).
D'après
Limits, lemmes \ref{limits-lemma-descend-finite-presentation} et
\ref{limits-lemma-descend-finite-locally-free}
il existe $i' \geq i$ et un morphisme fini localement libre
$Z_{i'} \to U_{i'}$ de degré $d$ dont le changement de base à $U$ est $Z$. D'après
Limits of Spaces, proposition
\ref{spaces-limits-proposition-characterize-locally-finite-presentation}
on peut, après avoir remplacé $i'$ par un indice plus grand, supposer qu'il existe un
morphisme $f_{i'} : Z_{i'} \to X_{y_i}$ tel que
$$
\xymatrix{
Z \ar[d] \ar[r] \ar@/^3ex/[rr]^{f_\xi} &
Z_{i'} \ar[d] \ar[r]_{f_{i'}} & X_{y_i} \ar[d] \\
U \ar[r] & U_{i'} \ar[r] & U_i
}
$$
soit commutatif. Posons
$\Xi_{i'} = (U_{i'}, Z_{i'}, y_{i'}, x_{i'}, \alpha_{i'})$
où
\begin{enumerate}
\item $y_{i'}$ est l'objet de $\mathcal{Y}$ au-dessus de $U_{i'}$
qui est l'image inverse de $y_i$ sur $U_{i'}$,
\item $x_{i'}$ est l'objet de $\mathcal{X}$ au-dessus de $Z_{i'}$ qui correspond,
par le lemme de $2$-Yoneda, au $1$-morphisme
$$
(\Sch/Z_{i'})_{fppf} \to
\mathcal{S}_{X_{y_i}} \to
(\Sch/U_i)_{fppf} \times_{y_i, \mathcal{Y}, F} \mathcal{X} \to
\mathcal{X}
$$
où la flèche du milieu est l'équivalence qui définit $X_{y_i}$
(notations comme dans
Algebraic Stacks, sections
\ref{algebraic-section-representable-by-algebraic-spaces} et
\ref{algebraic-section-split}).
\item $\alpha_{i'} : y_{i'}|_{Z_{i'}} \to F(x_{i'})$ est l'isomorphisme
provenant de la $2$-commutativité du diagramme
$$
\xymatrix{
(\Sch/Z_{i'})_{fppf} \ar[r] \ar[rd] &
(\Sch/U_i)_{fppf} \times_{y_i, \mathcal{Y}, F} \mathcal{X}
\ar[r] \ar[d] &
\mathcal{X} \ar[d]^F \\
& (\Sch/U_{i'})_{fppf} \ar[r] & \mathcal{Y}
}
$$
\end{enumerate}
Rappelons que $f_\xi : Z \to X_{y_i}$ était le morphisme correspondant à
l'objet $\xi = (Z, Z \to U_i, x, \alpha)$ de
$(\Sch/U_i)_{fppf} \times_{y_i, \mathcal{Y}, F} \mathcal{X}$
au-dessus de $Z$. Par construction, $f_{i'}$ est le morphisme correspondant à
l'objet $\xi_{i'} = (Z_{i'}, Z_{i'} \to U_i, x_{i'}, \alpha_{i'})$.
Comme $f_\xi = f_{i'} \circ (Z \to Z_{i'})$, on voit que
l'objet $\xi_{i'} = (Z_{i'}, Z_{i'} \to U_i, x_{i'}, \alpha_{i'})$ a pour image inverse
$\xi$ sur $Z$. Ainsi, $x_{i'}$ a pour image inverse $x$ et $\alpha_{i'}$
a pour image inverse $\alpha$. Cela signifie que $\Xi_{i'}$ a pour image inverse
$\Xi$ sur $U$, ce qui conclut.
\end{proof}










\section{Le champ de Hilbert fini d'un point}
\label{section-hilbert-point}

\noindent
Soit $d \geq 1$ un entier. Dans
Examples of Stacks, définition \ref{examples-stacks-definition-hilbert-d-stack},
nous avons défini un champ en groupoïdes $\mathcal{H}_d$.
Dans cette section, nous montrons que $\mathcal{H}_d$ est un
champ algébrique. Nous supposerons partout que
$S = \Spec(\mathbf{Z})$.
Le cas général s'en déduira par changement de base.
Rappelons que la catégorie fibre de $\mathcal{H}_d$ au-dessus d'un schéma $T$
est la catégorie des morphismes finis localement libres $\pi : Z \to T$ de
degré $d$. Au lieu de les classifier directement, nous étudions d'abord
les faisceaux quasi-cohérents d'algèbres $\pi_*\mathcal{O}_Z$.

\medskip\noindent
Soit $R$ un anneau. Adoptons temporairement la définition suivante :
une {\it algèbre libre de rang $d$ sur $R$}
est la donnée d'une structure $m$ de $R$-algèbre commutative sur $R^{\oplus d}$
dont $e_1 = (1, 0, \ldots, 0)$ est l'unité\footnote{Il est peut-être préférable
de la voir comme un couple formé d'une application de multiplication
$m : R^{\oplus d} \otimes_R R^{\oplus d} \to R^{\oplus d}$ et
d'un homomorphisme d'anneaux $\psi : R \to R^{\oplus d}$ satisfaisant à un certain nombre d'axiomes.}.
On considère $m$ comme une application $R$-linéaire
$$
m : R^{\oplus d} \otimes_R R^{\oplus d} \longrightarrow R^{\oplus d}
$$
telle que $m(e_1, x) = m(x, e_1) = x$ et que $m$ définisse une
structure d'anneau commutative et associative. Si l'on écrit
$m(e_i, e_j) = \sum a_{ij}^ke_k$, cela revient
aux conditions
$$
\left\{
\begin{matrix}
\sum_l a_{ij}^la_{lk}^m = \sum_l a_{il}^ma_{jk}^l & \forall i, j, k, m \\
a_{ij}^k = a_{ji}^k & \forall i, j, k \\
a_{i1}^j = \delta_{ij} & \forall i, j
\end{matrix}
\right.
$$
où $\delta_{ij}$ est le symbole de Kronecker. Définissons donc
$$
R_{univ} = \mathbf{Z}[a_{ij}^k]/J
$$
où $J$ est l'idéal engendré par les relations affichées ci-dessus.
Notons
$$
m_{univ} :
R_{univ}^{\oplus d} \otimes_{R_{univ}} R_{univ}^{\oplus d}
\longrightarrow
R_{univ}^{\oplus d}
$$
l'algèbre libre de rang $d$ sur $R_{univ}$ dont les constantes
de structure sont les classes des $a_{ij}^k$ modulo $J$.
Il est alors clair que, pour toute algèbre libre de rang $d$, notée $m$, sur un anneau
$R$, il existe un unique homomorphisme de $\mathbf{Z}$-algèbres
$\psi : R_{univ} \to R$ tel que $\psi_*m_{univ} = m$ (cela signifie que
$m$ s'obtient en appliquant le foncteur de changement de base
$- \otimes_{R_{univ}} R$ à $m_{univ}$). Autrement dit, en posant
$X = \Spec(R_{univ})$, on obtient une identification canonique
$$
X(T) = \{\text{algèbres libres de rang }d\text{ sur }R\}
$$
pour $T = \Spec(R)$ variable. Par localisation de Zariski, on obtient
l'identification suivante, apparemment plus générale,
\begin{equation}
\label{equation-objects}
X(T) = \{\text{algèbres libres de rang }d
\text{ sur }\Gamma(T, \mathcal{O}_T)\}
\end{equation}
pour tout schéma $T$.

\medskip\noindent
Parlons maintenant brièvement des {\it isomorphismes d'algèbres libres de rang $d$
sur $R$}. Supposons que $m$ et $m'$ soient deux algèbres libres de rang $d$
sur un anneau $R$. Un {\it isomorphisme de $m$ vers $m'$} est la donnée d'une
application $R$-linéaire inversible
$$
\varphi : R^{\oplus d} \longrightarrow R^{\oplus d}
$$
telle que $\varphi(e_1) = e_1$ et que
$$
m \circ \varphi \otimes \varphi = \varphi \circ m'.
$$
Notons que l'on peut composer ces isomorphismes, de sorte que la collection des
algèbres libres de rang $d$ sur $R$ devienne une catégorie.
On obtient ainsi un foncteur
\begin{equation}
\label{equation-FAd}
FA_d : \Sch_{fppf}^{opp} \longrightarrow \textit{Groupoids}
\end{equation}
de la catégorie des schémas vers les groupoïdes : à un schéma $T$, on associe
l'ensemble des algèbres libres de rang $d$ sur $\Gamma(T, \mathcal{O}_T)$,
muni d'une structure
de catégorie au moyen de la notion d'isomorphisme qui vient d'être définie.

\medskip\noindent
Ce qui précède suggère de considérer le foncteur en groupes $G$
qui associe à tout schéma $T$ le groupe
$$
G(T) = \{g \in \text{GL}_d(\Gamma(T, \mathcal{O}_T)) \mid g(e_1) = e_1\}
$$
Il est clair que $G \subset \text{GL}_d$ (voir
Groupoids, exemple \ref{groupoids-example-general-linear-group})
est le sous-schéma en groupes fermé défini par les équations
$x_{11} = 1$ et $x_{i1} = 0$ pour $i > 1$. Ainsi, $G$ est un
schéma en groupes affine et lisse sur $\Spec(\mathbf{Z})$. Considérons
l'action
$$
a : G \times_{\Spec(\mathbf{Z})} X \longrightarrow X
$$
qui associe à un point $(g, m)$ à valeurs dans $T$, où $T = \Spec(R)$,
du membre de gauche l'algèbre libre de rang $d$ sur $R$
donnée par
$$
a(g, m) = g^{-1} \circ m \circ g \otimes g.
$$
Notons que cela signifie que $g$ définit un isomorphisme $m \to a(g, m)$
d'algèbres libres de rang $d$ sur $R$. Nous omettons de vérifier que
$a$ définit bien une action du schéma en groupes $G$ sur le schéma $X$.

\begin{lemma}
\label{lemma-represent-FAd}
Le foncteur en groupoïdes $FA_d$ défini dans (\ref{equation-FAd})
est isomorphe (!) au foncteur en groupoïdes qui associe
à un schéma $T$ la catégorie dont
\begin{enumerate}
\item l'ensemble des objets est $X(T)$,
\item l'ensemble des morphismes est $G(T) \times X(T)$,
\item $s : G(T) \times X(T) \to X(T)$ est la projection,
\item $t : G(T) \times X(T) \to X(T)$ est $a(T)$, et
\item la composition $G(T) \times X(T) \times_{s, X(T), t} G(T) \times X(T)
\to G(T) \times X(T)$ est donnée par $((g, m), (g', m')) \mapsto (gg', m')$.
\end{enumerate}
\end{lemma}

\begin{proof}
Nous avons vu la règle sur les objets dans (\ref{equation-objects}).
Nous avons aussi vu ci-dessus que $g \in G(T)$ peut être considéré comme
un morphisme de $m$ vers $a(g, m)$ pour toute algèbre libre de rang $d$, notée $m$.
Réciproquement, tout morphisme $m \to m'$ est donné par une application linéaire
inversible $\varphi$ qui correspond à un élément $g \in G(T)$ tel
que $m' = a(g, m)$.
\end{proof}

\noindent
En fait, le groupoïde $(X, G \times X, s, t, c)$ décrit dans le
lemme ci-dessus est le groupoïde associé à l'action $a : G \times X \to X$
tel qu'il est défini dans
Groupoids, lemme \ref{groupoids-lemma-groupoid-from-action}.
Comme $G$ est lisse sur $\Spec(\mathbf{Z})$,
les deux morphismes $s, t : G \times X \to X$ sont
lisses : par symétrie, il suffit de le montrer pour l'un d'eux, et
$s$ est le changement de base de $G \to \Spec(\mathbf{Z})$.
Ainsi, $(X, G \times X, s, t, c)$ est un groupoïde lisse en schémas,
et le champ quotient $[X/G]$ est un champ algébrique d'après
Algebraic Stacks,
théorème \ref{algebraic-theorem-smooth-groupoid-gives-algebraic-stack}.

\begin{proposition}
\label{proposition-finite-hilbert-point}
Le champ $\mathcal{H}_d$ est équivalent au champ quotient
$[X/G]$ décrit ci-dessus. En particulier, $\mathcal{H}_d$ est un
champ algébrique.
\end{proposition}

\begin{proof}
Notons que, d'après
Groupoids in Spaces, définition
\ref{spaces-groupoids-definition-quotient-stack}
le champ quotient $[X/G]$ est le champ associé à la
catégorie fibrée en groupoïdes associée au « préfaisceau en groupoïdes »
qui associe à un schéma $T$ le groupoïde
$$
(X(T), G(T) \times X(T), s, t, c).
$$
Comme ce « préfaisceau en groupoïdes » est isomorphe à $FA_d$ d'après le
lemme \ref{lemma-represent-FAd},
il suffit de montrer que $\mathcal{H}_d$ est le champ associé
à (la catégorie fibrée en groupoïdes associée au
« préfaisceau en groupoïdes ») $FA_d$. Pour cela, définissons d'abord un
foncteur
$$
\Spec : FA_d \longrightarrow \mathcal{H}_d
$$
Rappelons que la catégorie fibre de $\mathcal{H}_d$ au-dessus d'un schéma $T$
est la catégorie des morphismes finis localement libres $Z \to T$ de degré $d$.
Étant donnés un schéma $T$ et une algèbre libre de rang $d$
sur $\Gamma(T, \mathcal{O}_T)$, notée $m$, on peut lui associer l'objet
$$
Z = \underline{\Spec}_T(\mathcal{A})
$$
de $\mathcal{H}_{d, T}$,
où $\mathcal{A} = \mathcal{O}_T^{\oplus d}$ est munie d'une
structure de $\mathcal{O}_T$-algèbre au moyen de $m$. De plus, si $m'$ est
une seconde algèbre libre de rang $d$ sur $\Gamma(T, \mathcal{O}_T)$
et si $\varphi : m \to m'$ est un isomorphisme entre elles, alors
l'application $\mathcal{O}_T$-linéaire induite
$\varphi : \mathcal{O}_T^{\oplus d} \to \mathcal{O}_T^{\oplus d}$
induit un isomorphisme
$$
\varphi : \mathcal{A}' \longrightarrow \mathcal{A}
$$
de $\mathcal{O}_T$-algèbres quasi-cohérentes. Par conséquent,
$$
\underline{\Spec}_T(\varphi) :
\underline{\Spec}_T(\mathcal{A})
\longrightarrow
\underline{\Spec}_T(\mathcal{A}')
$$
est un morphisme dans la catégorie fibre $\mathcal{H}_{d, T}$. Nous omettons de
vérifier que cette construction est compatible au changement de base ; on
obtient donc bien un foncteur $\Spec : FA_d \to \mathcal{H}_d$,
comme annoncé ci-dessus.

\medskip\noindent
Pour montrer que $\Spec : FA_d \to \mathcal{H}_d$ induit une équivalence
entre le champ associé à $FA_d$ et $\mathcal{H}_d$, il suffit de
vérifier que
\begin{enumerate}
\item $\mathit{Isom}(m, m') = \mathit{Isom}(\Spec(m), \Spec(m'))$
pour tous $m, m' \in FA_d(T)$.
\item pour tout schéma $T$ et tout objet $Z \to T$ de $\mathcal{H}_{d, T}$,
il existe un recouvrement $\{T_i \to T\}$ tel que $Z|_{T_i}$ soit
isomorphe à $\Spec(m)$ pour un certain $m \in FA_d(T_i)$.
\end{enumerate}
Voir
Stacks, lemme \ref{stacks-lemma-stackify-groupoids}.
La première assertion résulte de l'observation que tout isomorphisme
$$
\underline{\Spec}_T(\mathcal{A})
\longrightarrow
\underline{\Spec}_T(\mathcal{A}')
$$
est nécessairement donné par une matrice globale inversible $g$ lorsque
$\mathcal{A} = \mathcal{A}' = \mathcal{O}_T^{\oplus d}$ comme modules.
Pour démontrer la seconde assertion, soit $\pi : Z \to T$ un morphisme fini
localement libre de degré $d$. Alors $\mathcal{A}=\pi_*\mathcal{O}_Z$ est un faisceau
de $\mathcal{O}_T$-modules localement libre de rang $d$.
Considérons l'élément $1 \in \Gamma(T, \mathcal{A})$. Cet élément est
non nul dans $\mathcal{A}_t \otimes_{\mathcal{O}_{T, t}} \kappa(t)$
pour tout $t \in T$, puisque le schéma
$Z_t = \Spec(\mathcal{A}_t \otimes_{\mathcal{O}_{T, t}} \kappa(t))$
est non vide, étant de degré $d > 0$ sur $\kappa(t)$. Ainsi,
$1 : \mathcal{O}_T \to \mathcal{A}$ peut localement servir de premier
élément d'une base (on peut, par exemple, utiliser les parties (1) et (2) du
lemme \ref{algebra-lemma-cokernel-flat} de Algebra
pour le voir). Après localisation sur
$T$, on peut donc supposer qu'il existe un isomorphisme
$\varphi : \mathcal{A} \to \mathcal{O}_T^{\oplus d}$
tel que $1 \in \Gamma(\mathcal{A})$ corresponde au premier élément de la base.
Dans cette situation, l'application de multiplication
$\mathcal{A} \otimes_{\mathcal{O}_T} \mathcal{A} \to \mathcal{A}$
se traduit, par $\varphi$, en une algèbre libre de rang $d$, notée $m$, sur
$\Gamma(T, \mathcal{O}_T)$. Cela achève la preuve.
\end{proof}




\section{Champs de Hilbert finis d'espaces}
\label{section-spaces-hilbert}

\noindent
Le champ de Hilbert fini d'un espace algébrique est un champ algébrique.

\begin{lemma}
\label{lemma-hilbert-stack-of-space}
Soit $S$ un schéma.
Soit $X$ un espace algébrique sur $S$.
Alors $\mathcal{H}_d(X)$ est un champ algébrique.
\end{lemma}

\begin{proof}
Le $1$-morphisme
$$
\mathcal{H}_d(X) \longrightarrow \mathcal{H}_d
$$
est représentable par des espaces algébriques d'après le
lemme \ref{lemma-representable-on-top}.
Le champ $\mathcal{H}_d$ est algébrique d'après la
proposition \ref{proposition-finite-hilbert-point}.
Par conséquent, $\mathcal{H}_d(X)$ est un champ algébrique d'après
Algebraic Stacks,
lemme \ref{algebraic-lemma-representable-morphism-to-algebraic}.
\end{proof}

\noindent
Ce lemme permet d'amorcer le raisonnement.

\begin{lemma}
\label{lemma-hilbert-stack-relative-space}
Soit $S$ un schéma. Soit $F : \mathcal{X} \to \mathcal{Y}$ un $1$-morphisme
de champs en groupoïdes sur $(\Sch/S)_{fppf}$ tel que
\begin{enumerate}
\item $\mathcal{X}$ soit représentable par un espace algébrique, et
\item $F$ soit représentable par des espaces algébriques, surjectif, plat et
localement de présentation finie.
\end{enumerate}
Alors $\mathcal{H}_d(\mathcal{X}/\mathcal{Y})$ est un champ algébrique.
\end{lemma}

\begin{proof}
Choisissons un champ en groupoïdes représentable $\mathcal{U}$ sur $S$ et un
$1$-morphisme $f : \mathcal{U} \to \mathcal{H}_d(\mathcal{X})$
qui soit représentable par des espaces algébriques, lisse et surjectif.
C'est possible car $\mathcal{H}_d(\mathcal{X})$ est un champ algébrique d'après le
lemme \ref{lemma-hilbert-stack-of-space}.
Considérons le $2$-produit fibré
$$
\mathcal{W} =
\mathcal{H}_d(\mathcal{X}/\mathcal{Y})
\times_{\mathcal{H}_d(\mathcal{X}), f}
\mathcal{U}.
$$
Puisque $\mathcal{U}$ est représentable (c'est en particulier un champ en
groupoïdes discrets), il résulte de
Examples of Stacks, lemme \ref{examples-stacks-lemma-faithful-hilbert}
et de
Stacks, lemme \ref{stacks-lemma-2-fibre-product-gives-stack-in-setoids}
que $\mathcal{W}$ est un champ en groupoïdes discrets. Le $1$-morphisme
$\mathcal{W} \to \mathcal{H}_d(\mathcal{X}/\mathcal{Y})$ est
représentable par des espaces algébriques, lisse et surjectif comme changement
de base du morphisme $f$ (voir
Algebraic Stacks,
lemmes \ref{algebraic-lemma-base-change-representable-by-spaces} et
\ref{algebraic-lemma-base-change-representable-transformations-property}).
Ainsi, si l'on montre que $\mathcal{W}$ est représentable par un espace algébrique,
le lemme résultera de
Algebraic Stacks,
lemme \ref{algebraic-lemma-smooth-surjective-morphism-implies-algebraic}.

\medskip\noindent
La diagonale de $\mathcal{Y}$ est représentable par des espaces algébriques d'après le
lemme \ref{lemma-flat-finite-presentation-surjective-diagonal}.
On peut appliquer le
lemme \ref{lemma-relative-hilbert}
pour voir que le $1$-morphisme
$$
\mathcal{H}_d(\mathcal{X}/\mathcal{Y})
\longrightarrow
\mathcal{H}_d(\mathcal{X}) \times \mathcal{Y}
$$
est représentable par des espaces algébriques. Considérons le $2$-produit fibré
$$
\mathcal{V} =
\mathcal{H}_d(\mathcal{X}/\mathcal{Y})
\times_{(\mathcal{H}_d(\mathcal{X}) \times \mathcal{Y}), f \times F}
(\mathcal{U} \times \mathcal{X}).
$$
Le morphisme de projection $\mathcal{V} \to \mathcal{U} \times \mathcal{X}$
est représentable par des espaces algébriques comme changement de base du dernier
morphisme affiché. Par conséquent, $\mathcal{V}$ est un espace algébrique (voir
Bootstrap, lemme \ref{bootstrap-lemma-representable-by-spaces-over-space}
ou
Algebraic Stacks,
lemme \ref{algebraic-lemma-base-change-by-space-representable-by-space}).
Le $1$-morphisme $\mathcal{V} \to \mathcal{W}$ s'insère dans le diagramme
$2$-cartésien suivant
$$
\xymatrix{
\mathcal{V} \ar[d] \ar[r] & \mathcal{X} \ar[d]^F \\
\mathcal{W} \ar[r] & \mathcal{Y}
}
$$
car
$$
\mathcal{H}_d(\mathcal{X}/\mathcal{Y})
\times_{(\mathcal{H}_d(\mathcal{X}) \times \mathcal{Y}), f \times F}
(\mathcal{U} \times \mathcal{X})
=
(\mathcal{H}_d(\mathcal{X}/\mathcal{Y})
\times_{\mathcal{H}_d(\mathcal{X}), f}
\mathcal{U}) \times_{\mathcal{Y}, F} \mathcal{X}.
$$
Ainsi, $\mathcal{V} \to \mathcal{W}$ est représentable par des espaces algébriques,
surjectif, plat et localement de présentation finie comme changement de base
de $F$. Il en est donc de même pour le morphisme correspondant entre les faisceaux
d'ensembles associés à $\mathcal{V}$ et à $\mathcal{W}$, voir
Algebraic Stacks, lemme \ref{algebraic-lemma-map-fibred-setoids-property}.
On en conclut que le faisceau associé à $\mathcal{W}$ est un espace
algébrique d'après
Bootstrap, théorème \ref{bootstrap-theorem-final-bootstrap}.
\end{proof}




\section{Lieu LCI dans le champ de Hilbert}
\label{section-lci}

\noindent
On se reportera à
Examples of Stacks, section \ref{examples-stacks-section-hilbert-d-stack}
pour les notations. Fixons un $1$-morphisme $F : \mathcal{X} \longrightarrow \mathcal{Y}$
de champs en groupoïdes sur $(\Sch/S)_{fppf}$. Supposons que
$F$ soit représentable par des espaces algébriques. Fixons $d \geq 1$. Considérons un
objet $(U, Z, y, x, \alpha)$ de $\mathcal{H}_d(\mathcal{X}/\mathcal{Y})$.
On en déduit un $1$-morphisme
$$
(\Sch/Z)_{fppf}
\longrightarrow
(\Sch/U)_{fppf} \times_{y, \mathcal{Y}, F} \mathcal{X}
$$
(d'après la propriété universelle des $2$-produits fibrés) qui est représenté par
un morphisme d'espaces algébriques sur $U$.
En effet, puisque $F$ est représentable par des espaces algébriques, on peut choisir
un espace algébrique $X_y$ sur $U$ qui représente le $2$-produit fibré
$(\Sch/U)_{fppf} \times_{y, \mathcal{Y}, F} \mathcal{X}$.
Comme $\alpha : y|_Z \to F(x)$ est un isomorphisme, on voit que
$\xi = (Z, Z \to U, x, \alpha)$ est un objet du $2$-produit fibré
$(\Sch/U)_{fppf} \times_{y, \mathcal{Y}, F} \mathcal{X}$ au-dessus de $Z$.
Ainsi, $\xi$ donne naissance à un morphisme $x_\alpha : Z \to X_y$ d'espaces algébriques
sur $U$, puisque $X_y$ est le foncteur des classes d'isomorphisme d'objets de
$(\Sch/U)_{fppf} \times_{y, \mathcal{Y}, F} \mathcal{X}$, voir
Algebraic Stacks,
lemme \ref{algebraic-lemma-characterize-representable-by-space}.
Voici un diagramme
\begin{equation}
\label{equation-relative-map}
\vcenter{
\xymatrix{
Z \ar[r]_{x_\alpha} \ar[rd] & X_y \ar[d] \\
& U
}
}
\quad\quad
\vcenter{
\xymatrix{
(\Sch/Z)_{fppf} \ar[rd] \ar[r]_-{x, \alpha} &
(\Sch/U)_{fppf} \times_{y, \mathcal{Y}, F} \mathcal{X} \ar[r] \ar[d] &
\mathcal{X} \ar[d]^F \\
& (\Sch/U)_{fppf} \ar[r]^y & \mathcal{Y}
}
}
\end{equation}
Remarquons que si
$(f, g, b, a) : (U, Z, y, x, \alpha) \to (U', Z', y', x', \alpha')$
est un morphisme entre objets de $\mathcal{H}_d$, alors le morphisme
$x'_{\alpha'} : Z' \to X'_{y'}$ est le changement de base du morphisme
$x_\alpha$ par le morphisme $g : U' \to U$ (détails omis).

\medskip\noindent
Supposons maintenant, de plus, que $F$ soit plat et localement de présentation finie.
Dans cette situation, définissons une sous-catégorie pleine
$$
\mathcal{H}_{d, lci}(\mathcal{X}/\mathcal{Y}) \subset
\mathcal{H}_d(\mathcal{X}/\mathcal{Y})
$$
formée des objets $(U, Z, y, x, \alpha)$ de
$\mathcal{H}_d(\mathcal{X}/\mathcal{Y})$ tels
que le morphisme correspondant $x_\alpha : Z \to X_y$ soit non ramifié
et localement d'intersection complète (voir
Morphisms of Spaces, définition \ref{spaces-morphisms-definition-unramified}
et
More on Morphisms of Spaces,
définition \ref{spaces-more-morphisms-definition-lci}
pour les définitions correspondantes).

\begin{lemma}
\label{lemma-lci-locus-stack-in-groupoids}
Soit $S$ un schéma. Fixons un $1$-morphisme
$F : \mathcal{X} \longrightarrow \mathcal{Y}$
de champs en groupoïdes sur $(\Sch/S)_{fppf}$.
Supposons que $F$ soit représentable par des espaces algébriques, plat et localement
de présentation finie. Alors $\mathcal{H}_{d, lci}(\mathcal{X}/\mathcal{Y})$
est un champ en groupoïdes et le foncteur d'inclusion
$$
\mathcal{H}_{d, lci}(\mathcal{X}/\mathcal{Y})
\longrightarrow
\mathcal{H}_d(\mathcal{X}/\mathcal{Y})
$$
est représentable et est une immersion ouverte.
\end{lemma}

\begin{proof}
Soit $\Xi = (U, Z, y, x, \alpha)$ un objet de $\mathcal{H}_d$. Il résulte
de la remarque qui suit
(\ref{equation-relative-map})
que l'image inverse de $\Xi$ par $U' \to U$ appartient à
$\mathcal{H}_{d, lci}(\mathcal{X}/\mathcal{Y})$ si et seulement si le changement
de base de $x_\alpha$ est non ramifié et localement d'intersection complète.
Notons que $Z \to U$ est fini localement libre (donc plat, localement de
présentation finie et universellement fermé) et que $X_y \to U$ est
plat et localement de présentation finie par notre hypothèse sur $F$. Alors
More on Morphisms of Spaces, lemmes
\ref{spaces-more-morphisms-lemma-where-unramified} et
\ref{spaces-more-morphisms-lemma-where-lci}
entraînent qu'il existe un sous-schéma ouvert $W \subset U$ tel qu'un morphisme
$U' \to U$ se factorise par $W$ si et seulement si le changement de base de
$x_\alpha$ par $U' \to U$ est non ramifié et localement d'intersection
complète. Cela implique que
$$
(\Sch/U)_{fppf}
\times_{\Xi, \mathcal{H}_d(\mathcal{X}/\mathcal{Y})}
\mathcal{H}_{d, lci}(\mathcal{X}/\mathcal{Y})
$$
est représentable par $W$. La dernière assertion du lemme
en résulte. La première assertion (à savoir que
$\mathcal{H}_{d, lci}(\mathcal{X}/\mathcal{Y})$ est un champ en groupoïdes)
résulte alors de ce qui précède et de
Algebraic Stacks,
lemme \ref{algebraic-lemma-open-fibred-category-is-algebraic}.
\end{proof}

\noindent
Les morphismes localement d'intersection complète sont « localement sans obstruction ».
Cette assertion vaut dans une bien plus grande généralité que le cas particulier
dont nous avons besoin ici.

\begin{lemma}
\label{lemma-lci-unobstructed}
Soit $U \subset U'$ un épaississement du premier ordre de schémas affines.
Soit $X'$ un espace algébrique plat sur $U'$. Posons $X = U \times_{U'} X'$.
Soit $Z \to U$ fini localement libre de degré $d$. Enfin, soit
$f : Z \to X$ non ramifié et localement d'intersection complète.
Alors il existe un diagramme commutatif
$$
\xymatrix{
(Z \subset Z') \ar[rd] \ar[rr]_{(f, f')} & & (X \subset X') \ar[ld] \\
& (U \subset U')
}
$$
d'espaces algébriques sur $U'$ tel que $Z' \to U'$ soit fini localement libre
de degré $d$ et que $Z = U \times_{U'} Z'$.
\end{lemma}

\begin{proof}
D'après
More on Morphisms of Spaces,
lemme \ref{spaces-more-morphisms-lemma-unramified-lci},
le faisceau conormal $\mathcal{C}_{Z/X}$ du morphisme non ramifié $Z \to X$
est un $\mathcal{O}_Z$-module localement libre de type fini et, d'après
More on Morphisms of Spaces,
lemme \ref{spaces-more-morphisms-lemma-transitivity-conormal-lci},
on a une suite exacte
$$
0 \to i^*\mathcal{C}_{X/X'} \to
\mathcal{C}_{Z/X'} \to
\mathcal{C}_{Z/X} \to 0
$$
de faisceaux conormaux. Comme $Z$ est affine, cette suite est scindée. Choisissons
un scindage
$$
\mathcal{C}_{Z/X'} = i^*\mathcal{C}_{X/X'} \oplus \mathcal{C}_{Z/X}
$$
Soit $Z \subset Z''$ l'épaississement universel du premier ordre de $Z$
sur $X'$ (voir
More on Morphisms of Spaces,
section \ref{spaces-more-morphisms-section-universal-thickening}).
Notons $\mathcal{I} \subset \mathcal{O}_{Z''}$ le faisceau quasi-cohérent
d'idéaux correspondant à $Z \subset Z''$. Par définition,
$\mathcal{C}_{Z/X'}$ est $\mathcal{I}$ considéré comme faisceau sur $Z$.
Le scindage ci-dessus détermine donc un scindage
$$
\mathcal{I} = i^*\mathcal{C}_{X/X'} \oplus \mathcal{C}_{Z/X}
$$
Soit $Z' \subset Z''$ le sous-schéma fermé défini par
$\mathcal{C}_{Z/X} \subset \mathcal{I}$, considéré comme faisceau quasi-cohérent
d'idéaux sur $Z''$. Il est clair que $Z'$ est un épaississement du premier ordre
de $Z$ et que l'on obtient un diagramme commutatif d'épaississements du premier ordre
comme dans l'énoncé du lemme.

\medskip\noindent
Puisque $X' \to U'$ est plat et que $X = U \times_{U'} X'$, on voit que
$\mathcal{C}_{X/X'}$ est l'image inverse de $\mathcal{C}_{U/U'}$ sur $X$, voir
More on Morphisms of Spaces, lemme \ref{spaces-more-morphisms-lemma-deform}.
Notons que, par construction, $\mathcal{C}_{Z/Z'} = i^*\mathcal{C}_{X/X'}$ ;
on en conclut que $\mathcal{C}_{Z/Z'}$ est isomorphe à l'image inverse
de $\mathcal{C}_{U/U'}$ sur $Z$. En appliquant
More on Morphisms of Spaces, lemme \ref{spaces-more-morphisms-lemma-deform}
une nouvelle fois (ou son analogue pour les schémas, voir
More on Morphisms, lemme \ref{more-morphisms-lemma-deform}),
on en conclut que $Z' \to U'$ est plat et que $Z = U \times_{U'} Z'$.
Enfin,
More on Morphisms, lemme \ref{more-morphisms-lemma-deform-property}
montre que $Z' \to U'$ est fini localement libre de degré $d$.
\end{proof}

\begin{lemma}
\label{lemma-lci-formally-smooth}
Soit $F : \mathcal{X} \to \mathcal{Y}$ un $1$-morphisme de champs en groupoïdes
sur $(\Sch/S)_{fppf}$. Supposons que $F$ soit représentable par des espaces
algébriques, plat et localement de présentation finie. Alors
$$
p : \mathcal{H}_{d, lci}(\mathcal{X}/\mathcal{Y}) \to \mathcal{Y}
$$
est formellement lisse sur les objets.
\end{lemma}

\begin{proof}
Il s'agit de montrer ceci : étant donnés
\begin{enumerate}
\item un objet $(U, Z, y, x, \alpha)$ de
$\mathcal{H}_{d, lci}(\mathcal{X}/\mathcal{Y})$ au-dessus d'un schéma affine $U$,
\item un épaississement du premier ordre $U \subset U'$, et
\item un objet $y'$ de $\mathcal{Y}$ au-dessus de $U'$ tel que $y'|_U = y$,
\end{enumerate}
il existe alors un objet $(U', Z', y', x', \alpha')$ de
$\mathcal{H}_{d, lci}(\mathcal{X}/\mathcal{Y})$ au-dessus de $U'$ tel que
$Z = U \times_{U'} Z'$, que $x = x'|_Z$ et que
$\alpha = \alpha'|_U$. En effet, les deux dernières égalités assureront
la commutativité de (\ref{equation-formally-smooth}).

\medskip\noindent
Considérons le morphisme $x_\alpha : Z \to X_y$ construit dans
l'équation (\ref{equation-relative-map}). Notons de même $X'_{y'}$
l'espace algébrique sur $U'$ qui représente le $2$-produit fibré
$(\Sch/U')_{fppf} \times_{y', \mathcal{Y}, F} \mathcal{X}$.
Par hypothèse, le morphisme $X'_{y'} \to U'$ est plat (et localement de
présentation finie). Comme $y'|_U = y$, on voit que $X_y = U \times_{U'} X'_{y'}$.
On peut donc appliquer le
lemme \ref{lemma-lci-unobstructed}
pour trouver $Z' \to U'$ fini localement libre de degré $d$ tel que
$Z = U \times_{U'} Z'$ et tel que $Z' \to X'_{y'}$ prolonge $x_\alpha$.
Par construction, le morphisme $Z' \to X'_{y'}$ correspond à un couple
$(x', \alpha')$. Il est clair que $(U', Z', y', x', \alpha')$
est un objet de $\mathcal{H}_d(\mathcal{X}/\mathcal{Y})$ au-dessus de $U'$
tel que $Z = U \times_{U'} Z'$, que $x = x'|_Z$ et que
$\alpha = \alpha'|_U$. Comme le
lemme \ref{lemma-lci-locus-stack-in-groupoids}
montre que $\mathcal{H}_{d, lci}(\mathcal{X}/\mathcal{Y}) \subset
\mathcal{H}_d(\mathcal{X}/\mathcal{Y})$ est un « sous-champ ouvert »,
il s'ensuit que $(U', Z', y', x', \alpha')$ est un objet de
$\mathcal{H}_{d, lci}(\mathcal{X}/\mathcal{Y})$, comme voulu.
\end{proof}

\begin{lemma}
\label{lemma-lci-surjective}
Soit $F : \mathcal{X} \to \mathcal{Y}$ un $1$-morphisme de champs en groupoïdes
sur $(\Sch/S)_{fppf}$. Supposons que $F$ soit représentable par des espaces
algébriques, plat, surjectif et localement de présentation finie. Alors
$$
\coprod\nolimits_{d \geq 1} \mathcal{H}_{d, lci}(\mathcal{X}/\mathcal{Y})
\longrightarrow
\mathcal{Y}
$$
est surjectif sur les objets.
\end{lemma}

\begin{proof}
Il suffit de montrer ceci : pour tout corps $k$
et tout objet $y$ de $\mathcal{Y}$ au-dessus de $\Spec(k)$, il existe
un entier $d \geq 1$ et un objet $(U, Z, y, x, \alpha)$ de
$\mathcal{H}_{d, lci}(\mathcal{X}/\mathcal{Y})$ avec $U = \Spec(k)$.
En effet, dans ce cas, on voit que $p$ est surjectif sur les objets au
sens fort où aucune extension du corps n'est nécessaire.

\medskip\noindent
Notons $X_y$ l'espace algébrique sur $U = \Spec(k)$
qui représente le $2$-produit fibré
$(\Sch/U)_{fppf} \times_{y, \mathcal{Y}, F} \mathcal{X}$.
Par hypothèse, le morphisme $X_y \to \Spec(k)$ est surjectif et
localement de présentation finie (et plat). En particulier, $X_y$ est
non vide. Choisissons un schéma affine non vide $V$ et un morphisme étale
$V \to X_y$. Notons que $V \to \Spec(k)$ est plat, surjectif
et localement de présentation finie (d'après
Morphisms of Spaces,
définition \ref{spaces-morphisms-definition-locally-finite-presentation}).
Choisissons un point fermé $v \in V$ où $V \to \Spec(k)$ est Cohen-Macaulay
(c'est-à-dire que $V$ est Cohen-Macaulay en $v$), voir
More on Morphisms,
lemme \ref{more-morphisms-lemma-flat-finite-presentation-CM-open}.
En appliquant
More on Morphisms,
lemme \ref{more-morphisms-lemma-slice-CM},
on trouve une immersion régulière $Z \to V$ telle que $Z = \{v\}$.
Cela implique que $Z \to V$ est une immersion fermée. De plus, il s'ensuit que
$Z \to \Spec(k)$ est fini (par exemple d'après
Algebra, lemme \ref{algebra-lemma-isolated-point}).
Ainsi, $Z \to \Spec(k)$ est fini localement libre d'un certain degré $d$.
Le morphisme $Z \to X_y$ est alors non ramifié comme composé
d'une immersion fermée et d'un morphisme étale
(voir
Morphisms of Spaces, lemmes \ref{spaces-morphisms-lemma-composition-unramified},
\ref{spaces-morphisms-lemma-etale-unramified} et
\ref{spaces-morphisms-lemma-immersion-unramified}).
Enfin, $Z \to X_y$ est localement d'intersection complète
comme composé d'une immersion régulière de schémas et d'un morphisme
étale d'espaces algébriques (voir
More on Morphisms, lemme \ref{more-morphisms-lemma-regular-immersion-lci}
et
Morphisms of Spaces, lemmes \ref{spaces-morphisms-lemma-etale-smooth} et
\ref{spaces-morphisms-lemma-smooth-syntomic}, ainsi que
More on Morphisms of Spaces,
lemmes \ref{spaces-more-morphisms-lemma-flat-lci} et
\ref{spaces-more-morphisms-lemma-composition-lci}).
Le morphisme $Z \to X_y$ correspond à un objet $x$ de $\mathcal{X}$
au-dessus de $Z$, muni d'un isomorphisme $\alpha : y|_Z \to F(x)$.
On obtient un objet $(U, Z, y, x, \alpha)$ de
$\mathcal{H}_d(\mathcal{X}/\mathcal{Y})$. D'après ce qui précède au sujet
du morphisme $Z \to X_y$, on voit qu'il s'agit en fait d'un objet de la
sous-catégorie $\mathcal{H}_{d, lci}(\mathcal{X}/\mathcal{Y})$, ce qui conclut.
\end{proof}




















\section{Amorçage des champs algébriques}
\label{section-bootstrap}

\noindent
Le théorème suivant est l'un des principaux résultats de ce chapitre.

\begin{theorem}
\label{theorem-bootstrap}
\begin{slogan}
Théorème d'Artin sur la représentabilité des groupoïdes plats
\end{slogan}
Soit $S$ un schéma. Soit $F : \mathcal{X} \to \mathcal{Y}$
un $1$-morphisme de champs en groupoïdes sur $(\Sch/S)_{fppf}$. Si
\begin{enumerate}
\item $\mathcal{X}$ est représentable par un espace algébrique, et
\item $F$ est représentable par des espaces algébriques, surjectif, plat et
localement de présentation finie,
\end{enumerate}
alors $\mathcal{Y}$ est un champ algébrique.
\end{theorem}

\begin{proof}
D'après le
lemme \ref{lemma-flat-finite-presentation-surjective-diagonal},
on voit que la diagonale de $\mathcal{Y}$ est représentable par des espaces
algébriques. Il suffit donc de vérifier l'existence d'un $1$-morphisme
$f : \mathcal{V} \to \mathcal{Y}$ de champs en groupoïdes sur
$(\Sch/S)_{fppf}$, avec $\mathcal{V}$ représentable et
$f$ surjectif et lisse. D'après le
lemme \ref{lemma-hilbert-stack-relative-space},
on sait que
$$
\coprod\nolimits_{d \geq 1} \mathcal{H}_d(\mathcal{X}/\mathcal{Y})
$$
est un champ algébrique. Il résulte du
lemme \ref{lemma-lci-locus-stack-in-groupoids}
et de
Algebraic Stacks,
lemme \ref{algebraic-lemma-open-fibred-category-is-algebraic}
que
$$
\coprod\nolimits_{d \geq 1} \mathcal{H}_{d, lci}(\mathcal{X}/\mathcal{Y})
$$
est lui aussi un champ algébrique. Choisissons un champ en groupoïdes représentable
$\mathcal{V}$ sur $(\Sch/S)_{fppf}$ et un $1$-morphisme surjectif et
lisse
$$
\mathcal{V}
\longrightarrow
\coprod\nolimits_{d \geq 1} \mathcal{H}_{d, lci}(\mathcal{X}/\mathcal{Y}).
$$
Nous affirmons que le composé
$$
\mathcal{V}
\longrightarrow
\coprod\nolimits_{d \geq 1} \mathcal{H}_{d, lci}(\mathcal{X}/\mathcal{Y})
\longrightarrow
\mathcal{Y}
$$
est lisse et surjectif, ce qui achève la preuve du théorème. En effet,
la lissité résultera des
lemmes \ref{lemma-limit-preserving} et \ref{lemma-lci-formally-smooth},
et la surjectivité du
lemme \ref{lemma-lci-surjective}.
Détaillons cela dans le paragraphe suivant.

\medskip\noindent
Par construction, $\mathcal{V} \to
\coprod\nolimits_{d \geq 1} \mathcal{H}_{d, lci}(\mathcal{X}/\mathcal{Y})$
est représentable par des espaces algébriques, surjectif et lisse (donc
aussi localement de présentation finie et formellement lisse d'après le principe
général
Algebraic Stacks, lemme
\ref{algebraic-lemma-representable-transformations-property-implication}
et
More on Morphisms of Spaces,
lemme \ref{spaces-more-morphisms-lemma-smooth-formally-smooth}).
En appliquant les
lemmes \ref{lemma-representable-by-spaces-limit-preserving},
\ref{lemma-representable-by-spaces-formally-smooth} et
\ref{lemma-representable-by-spaces-surjective}
on voit que $\mathcal{V} \to
\coprod\nolimits_{d \geq 1} \mathcal{H}_{d, lci}(\mathcal{X}/\mathcal{Y})$
commute aux limites sur les objets, est formellement lisse sur les objets et
est surjectif sur les objets. Le $1$-morphisme
$\coprod\nolimits_{d \geq 1} \mathcal{H}_{d, lci}(\mathcal{X}/\mathcal{Y})
\to \mathcal{Y}$ possède les propriétés suivantes :
\begin{enumerate}
\item commute aux limites sur les objets : cela résulte du
lemme \ref{lemma-limit-preserving}
pour $\mathcal{H}_d(\mathcal{X}/\mathcal{Y}) \to \mathcal{Y}$,
que l'on combine aux lemmes
\ref{lemma-lci-locus-stack-in-groupoids},
\ref{lemma-open-immersion-limit-preserving} et
\ref{lemma-composition-limit-preserving}
pour l'obtenir pour $\mathcal{H}_{d, lci}(\mathcal{X}/\mathcal{Y}) \to \mathcal{Y}$ ;
\item est formellement lisse sur les objets d'après le
lemme \ref{lemma-lci-formally-smooth} ;

\item est surjectif sur les objets d'après le
lemme \ref{lemma-lci-surjective}.
\end{enumerate}
En utilisant les
lemmes \ref{lemma-composition-limit-preserving},
\ref{lemma-composition-formally-smooth} et
\ref{lemma-composition-surjective}
on en conclut que le composé $\mathcal{V} \to \mathcal{Y}$
commute aux limites sur les objets, est formellement lisse sur les objets et
est surjectif sur les objets.
En utilisant les
lemmes \ref{lemma-representable-by-spaces-limit-preserving},
\ref{lemma-representable-by-spaces-formally-smooth} et
\ref{lemma-representable-by-spaces-surjective}
on voit que $\mathcal{V} \to \mathcal{Y}$ est
localement de présentation finie, formellement lisse et surjectif.
Enfin, en utilisant (par l'intermédiaire du principe général
Algebraic Stacks,
lemme \ref{algebraic-lemma-representable-transformations-property-implication})
le critère infinitésimal de relèvement
(More on Morphisms of Spaces, lemme
\ref{spaces-more-morphisms-lemma-smooth-formally-smooth})
on voit que $\mathcal{V} \to \mathcal{Y}$ est lisse, ce qui conclut.
\end{proof}








\section{Applications}
\label{section-applications}

\noindent
Notre première tâche consiste à montrer que le champ quotient $[U/R]$ associé à
un « groupoïde plat et localement de présentation finie » est un champ algébrique.
Voir
Groupoids in Spaces,
définition \ref{spaces-groupoids-definition-quotient-stack}
pour la définition du champ quotient.
Le lemme préliminaire suivant est l'analogue de
Algebraic Stacks,
lemme \ref{algebraic-lemma-smooth-quotient-smooth-presentation}.

\begin{lemma}
\label{lemma-flat-quotient-flat-presentation}
Soit $S$ un schéma appartenant à $\Sch_{fppf}$.
Soit $(U, R, s, t, c)$ un groupoïde en espaces algébriques sur $S$.
Supposons que $s$ et $t$ soient plats et localement de présentation finie.
Alors le morphisme $\mathcal{S}_U \to [U/R]$ est plat, localement de
présentation finie et surjectif.
\end{lemma}

\begin{proof}
Soit $T$ un schéma et soit $x : (\Sch/T)_{fppf} \to [U/R]$
un $1$-morphisme. Il faut montrer que la projection
$$
\mathcal{S}_U \times_{[U/R]} (\Sch/T)_{fppf}
\longrightarrow
(\Sch/T)_{fppf}
$$
est surjective, plate et localement de présentation finie.
On sait déjà que le membre de gauche
est représentable par un espace algébrique $F$, voir
Algebraic Stacks, lemmes \ref{algebraic-lemma-diagonal-quotient-stack} et
\ref{algebraic-lemma-representable-diagonal}.
Il faut donc montrer que le morphisme correspondant $F \to T$ d'espaces
algébriques est surjectif, localement de présentation finie et plat.
Puisqu'il s'agit de propriétés de morphismes d'espaces
algébriques qui sont locales sur le but pour la topologie fppf, on
peut les vérifier localement pour la topologie fppf sur $T$. Par construction, il existe
un recouvrement fppf $\{T_i \to T\}$ de $T$ tel que
$x|_{(\Sch/T_i)_{fppf}}$ provienne d'un morphisme
$x_i : T_i \to U$. (Notons que $F \times_T T_i$ représente le
$2$-produit fibré $\mathcal{S}_U \times_{[U/R]} (\Sch/T_i)_{fppf}$,
de sorte que tout est compatible au changement de base par $T_i \to T$.)
On peut donc supposer que $x$ provient de $x : T \to U$.
Dans ce cas, on voit que
$$
\mathcal{S}_U \times_{[U/R]} (\Sch/T)_{fppf}
=
(\mathcal{S}_U \times_{[U/R]} \mathcal{S}_U)
\times_{\mathcal{S}_U} (\Sch/T)_{fppf}
=
\mathcal{S}_R \times_{\mathcal{S}_U} (\Sch/T)_{fppf}
$$
La première égalité résulte de
Categories, lemme \ref{categories-lemma-2-fibre-product-erase-factor},
et la seconde de
Groupoids in Spaces,
lemme \ref{spaces-groupoids-lemma-quotient-stack-2-cartesian}.
Manifestement, le dernier $2$-produit fibré est représenté par l'espace
algébrique $F = R \times_{s, U, x} T$, et la projection
$R \times_{s, U, x} T \to T$ est plate et localement de présentation finie
comme changement de base du morphisme d'espaces algébriques
plat et localement de présentation finie $s : R \to U$.
Elle est aussi surjective puisque $s$ admet une section (à savoir la section unité
$e : U \to R$ du groupoïde).
Cela démontre le lemme.
\end{proof}

\noindent
Voici le premier résultat principal de cette section.

\begin{theorem}
\label{theorem-flat-groupoid-gives-algebraic-stack}
Soit $S$ un schéma appartenant à $\Sch_{fppf}$.
Soit $(U, R, s, t, c)$ un groupoïde en espaces algébriques sur $S$.
Supposons que $s$ et $t$ soient plats et localement de présentation finie.
Alors le champ quotient $[U/R]$ est un champ algébrique sur $S$.
\end{theorem}

\begin{proof}
Vérifions les deux conditions du
théorème \ref{theorem-bootstrap}
pour le morphisme
$$
(\Sch/U)_{fppf} \longrightarrow [U/R].
$$
La première est immédiate (car $U$ est un espace algébrique).
La seconde est le
lemme \ref{lemma-flat-quotient-flat-presentation}.
\end{proof}









\section{Quand un champ quotient est-il algébrique ?}
\label{section-quotient-algebraic}

\noindent
Dans
Groupoids in Spaces, section \ref{spaces-groupoids-section-stacks},
nous avons défini le champ quotient $[U/R]$ associé à un groupoïde
$(U, R, s, t, c)$ en espaces algébriques. Notons que $[U/R]$ est un champ
en groupoïdes dont la diagonale est représentable par des espaces algébriques (voir
Bootstrap, lemme \ref{bootstrap-lemma-quotient-stack-isom}
et
Algebraic Stacks, lemme \ref{algebraic-lemma-representable-diagonal})
et qu'il existe un espace algébrique $U$ et un $1$-morphisme
$(\Sch/U)_{fppf} \to [U/R]$ qui est une « surjection fppf »
au sens où il induit un morphisme entre les préfaisceaux des classes d'isomorphisme
d'objets qui devient surjectif après passage aux faisceaux associés.
Cependant, $[U/R]$ n'est pas en général un champ
algébrique. Cela ne contredit pas le
théorème \ref{theorem-bootstrap},
car le $1$-morphisme $(\Sch/U)_{fppf} \to [U/R]$ peut ne pas
être plat et localement de présentation finie.

\medskip\noindent
La manière la plus simple de construire des exemples de champs quotients non algébriques
consiste à considérer des quotients de la forme $[S/G]$, où $S$ est un schéma et $G$
un schéma en groupes sur $S$ qui agit trivialement sur $S$. En effet, nous verrons
ci-dessous
(lemme \ref{lemma-BG-algebraic})
que, si $[S/G]$ est algébrique, alors $G \to S$ doit être plat et localement
de présentation finie. On trouvera un exemple explicite dans
Examples, section \ref{examples-section-not-algebraic-stack}.

\begin{lemma}
\label{lemma-quotient-algebraic}
Soit $S$ un schéma et soit $B$ un espace algébrique sur $S$.
Soit $(U, R, s, t, c)$ un groupoïde en espaces algébriques sur $B$.
Le champ quotient $[U/R]$ est un champ algébrique si et seulement s'il
existe un morphisme d'espaces algébriques $g : U' \to U$ tel que
\begin{enumerate}
\item le composé
$U' \times_{g, U, t} R \to R \xrightarrow{s} U$ est un morphisme surjectif de
faisceaux, et
\item les morphismes $s', t' : R' \to U'$ sont plats et localement de présentation
finie, où $(U', R', s', t', c')$ est la restriction de
$(U, R, s, t, c)$ par $g$.
\end{enumerate}
\end{lemma}

\begin{proof}
Supposons d'abord que $g : U' \to U$ satisfasse à (1) et (2). La propriété (1)
implique que $[U'/R'] \to [U/R]$ est une équivalence, voir
Groupoids in Spaces,
lemme \ref{spaces-groupoids-lemma-quotient-stack-restrict-equivalence}.
D'après le
théorème \ref{theorem-flat-groupoid-gives-algebraic-stack},
le champ quotient $[U'/R']$ est un champ algébrique. Par conséquent,
$[U/R]$ est lui aussi un champ algébrique, voir
Algebraic Stacks, lemme \ref{algebraic-lemma-equivalent}.

\medskip\noindent
Réciproquement, supposons que $[U/R]$ soit un champ algébrique. On peut choisir un
schéma $W$ et un $1$-morphisme lisse et surjectif
$$
f : (\Sch/W)_{fppf} \longrightarrow [U/R].
$$
D'après le lemme de $2$-Yoneda
(Algebraic Stacks, section \ref{algebraic-section-2-yoneda}),
cela correspond à un objet $\xi$ de $[U/R]$ au-dessus de $W$.
D'après la description de $[U/R]$ dans
Groupoids in Spaces, lemme \ref{spaces-groupoids-lemma-quotient-stack-objects},
on peut trouver un morphisme de schémas surjectif, plat et localement de présentation finie
$b : U' \to W$ tel que $\xi' = b^*\xi$ corresponde à un morphisme
$g : U' \to U$. Notons que le $1$-morphisme
$$
f' : (\Sch/U')_{fppf} \longrightarrow [U/R]
$$
correspondant à $\xi'$ est surjectif, plat et localement de
présentation finie, voir
Algebraic Stacks, lemme
\ref{algebraic-lemma-composition-representable-transformations-property}.
Par conséquent,
$(\Sch/U')_{fppf} \times_{[U/R]} (\Sch/U')_{fppf}$
qui est représenté par l'espace algébrique
$$
\mathit{Isom}_{[U/R]}(\text{pr}_0^*\xi', \text{pr}_1^*\xi') =
(U' \times_S U')
\times_{(g \circ \text{pr}_0, g \circ \text{pr}_1), U \times_S U} R = R'
$$
(voir
Groupoids in Spaces, lemme
\ref{spaces-groupoids-lemma-quotient-stack-morphisms}
pour la première égalité ; la seconde est la définition de la restriction)
est plat et localement de présentation finie sur $U'$ à la fois par $s'$ et par $t'$
(par changement de base, voir
Algebraic Stacks, lemme
\ref{algebraic-lemma-base-change-representable-transformations-property}).
D'après cette description de $R'$ et
Algebraic Stacks, lemme \ref{algebraic-lemma-map-space-into-stack},
on obtient un $1$-morphisme canonique pleinement fidèle $[U'/R'] \to [U/R]$.
Ce $1$-morphisme est essentiellement surjectif parce que $f'$ est plat,
localement de présentation finie et surjectif (voir
Stacks, lemme \ref{stacks-lemma-characterize-essentially-surjective-when-ff}) ;
une autre manière de le démontrer consiste à utiliser
Algebraic Stacks, remarque
\ref{algebraic-remark-flat-fp-presentation}.
Enfin, on peut utiliser
Groupoids in Spaces, lemme
\ref{spaces-groupoids-lemma-quotient-stack-restrict-equivalence}
pour conclure que le composé
$U' \times_{g, U, t} R \to R \xrightarrow{s} U$ est un morphisme surjectif de faisceaux.
\end{proof}

\begin{lemma}
\label{lemma-group-quotient-algebraic}
Soit $S$ un schéma et soit $B$ un espace algébrique sur $S$.
Soit $G$ un espace algébrique en groupes sur $B$.
Soit $X$ un espace algébrique sur $B$ et soit $a : G \times_B X \to X$
une action de $G$ sur $X$ au-dessus de $B$.
Le champ quotient $[X/G]$ est un champ algébrique si et seulement s'il
existe un morphisme d'espaces algébriques $\varphi : X' \to X$ tel que
\begin{enumerate}
\item $G \times_B X' \to X$, $(g, x') \mapsto a(g, \varphi(x'))$, soit un
morphisme surjectif de faisceaux, et
\item les deux projections $X'' \to X'$ de l'espace algébrique $X''$
défini par la règle
$$
T \longmapsto \{(x'_1, g, x'_2) \in (X' \times_B G \times_B X')(T)
\mid \varphi(x'_1) = a(g, \varphi(x'_2))\}
$$
soient plates et localement de présentation finie.
\end{enumerate}
\end{lemma}

\begin{proof}
Ce lemme est un cas particulier du
lemme \ref{lemma-quotient-algebraic}.
En effet, d'après
Groupoids in Spaces, définition \ref{spaces-groupoids-definition-quotient-stack},
le champ quotient $[X/G]$ est égal au champ quotient $[X/G \times_B X]$ du groupoïde
en espaces algébriques $(X, G \times_B X, s, t, c)$ associé à
l'action du groupe dans
Groupoids in Spaces, lemme \ref{spaces-groupoids-lemma-groupoid-from-action}.
Une petite observation est nécessaire pour obtenir la condition (1).
En effet, le morphisme $s : G \times_B X \to X$ est la seconde projection
et le morphisme $t :  G \times_B X \to X$ est le morphisme d'action $a$.
Ainsi, le morphisme $h : U' \times_{g, U, t} R \to R \xrightarrow{s} U$ du
lemme \ref{lemma-quotient-algebraic}
correspond au morphisme
$$
X' \times_{\varphi, X, a} (G \times_B X) \xrightarrow{\text{pr}_1} X
$$
dans la situation présente. Toutefois, grâce à la symétrie fournie par
le passage à l'inverse dans $G$, ce morphisme est isomorphe au morphisme
$$
(G \times_B X) \times_{\text{pr}_1, X, \varphi} X' \xrightarrow{a} X
$$
de l'énoncé du lemme. Les détails sont omis.
\end{proof}

\begin{lemma}
\label{lemma-BG-algebraic}
\begin{slogan}
Les gerbes sont algébriques si et seulement si les groupes associés sont plats
et localement de présentation finie
\end{slogan}
Soit $S$ un schéma et soit $B$ un espace algébrique sur $S$.
Soit $G$ un espace algébrique en groupes sur $B$.
Munissons $B$ de l'action triviale de $G$.
Alors le champ quotient $[B/G]$ est un champ algébrique
si et seulement si $G$ est plat et localement de présentation finie sur $B$.
\end{lemma}

\begin{proof}
Si $G$ est plat et localement de présentation finie sur $B$, alors
$[B/G]$ est un champ algébrique d'après le
théorème \ref{theorem-flat-groupoid-gives-algebraic-stack}.

\medskip\noindent
Réciproquement, supposons que $[B/G]$ soit un champ algébrique. D'après le
lemme \ref{lemma-group-quotient-algebraic}
et puisque l'action est triviale, on voit
qu'il existe un espace algébrique $B'$ et un morphisme
$B' \to B$ tels que (1) $B' \to B$ soit un morphisme surjectif
de faisceaux et (2) que les projections
$$
B' \times_B G \times_B B' \to B'
$$
soient plates et localement de présentation finie. Notons que le changement de base
$B' \times_B G \times_B B' \to G \times_B B'$ de $B' \to B$
est lui aussi un morphisme surjectif de faisceaux. Il résulte donc de
Descent on Spaces, lemme \ref{spaces-descent-lemma-curiosity},
que la projection $G \times_B B' \to B'$ est plate et localement
de présentation finie. D'après (1), on peut trouver un recouvrement fppf
$\{B_i \to B\}$ tel que $B_i \to B$ se factorise par $B' \to B$.
Ainsi, $G \times_B B_i \to B_i$ est plat et localement de présentation finie
par changement de base. D'après
Descent on Spaces, lemmes
\ref{spaces-descent-lemma-descending-property-flat} et
\ref{spaces-descent-lemma-descending-property-locally-finite-presentation}
on en conclut que $G \to B$ est plat et localement de présentation finie.
\end{proof}

\noindent
Nous verrons plus loin que le champ quotient d'un $S$-espace lisse
par un espace algébrique en groupes $G$ est lisse, même lorsque $G$ ne l'est pas
(Morphisms of Stacks, lemme
\ref{stacks-morphisms-lemma-smooth-quotient-stack}).




\section{Champs algébriques pour la topologie étale}
\label{section-stacks-etale}

\noindent
Soit $S$ un schéma. Au lieu de travailler avec des champs en groupoïdes sur
le grand site fppf $(\Sch/S)_{fppf}$, on pourrait travailler avec des champs en groupoïdes
sur le grand site étale $(\Sch/S)_\etale$. Tout le contenu de
Algebraic Stacks, sections
\ref{algebraic-section-representable},
\ref{algebraic-section-2-yoneda},
\ref{algebraic-section-representable-morphism},
\ref{algebraic-section-split},
\ref{algebraic-section-representable-by-algebraic-spaces},
\ref{algebraic-section-morphisms-representable-by-algebraic-spaces},
\ref{algebraic-section-representable-properties} et
\ref{algebraic-section-stacks}
s'applique aux catégories fibrées en groupoïdes sur $(\Sch/S)_\etale$.
On obtient ainsi une seconde notion de champ algébrique en travaillant pour la
topologie étale. Cette notion est a priori plus faible que celle introduite
dans Algebraic Stacks, définition \ref{algebraic-definition-algebraic-stack},
puisqu'un champ pour la topologie fppf est certainement un champ pour la topologie
étale. Toutefois, les deux notions sont équivalentes, comme le montre le
lemme suivant.

\begin{lemma}
\label{lemma-stacks-etale}
Notons $\Sch_\alpha$ la catégorie sous-jacente commune à $\Sch_{fppf}$
et à $\Sch_\etale$ (voir
Sheaves on Stacks, section \ref{stacks-sheaves-section-sheaves} et
Topologies, remarque \ref{topologies-remark-choice-sites}). Soit $S$ un objet
de $\Sch_\alpha$. Soit
$$
p : \mathcal{X} \to \Sch_\alpha/S
$$
une catégorie fibrée en groupoïdes ayant les propriétés suivantes :
\begin{enumerate}
\item $\mathcal{X}$ est un champ en groupoïdes sur $(\Sch/S)_\etale$,
\item la diagonale $\Delta : \mathcal{X} \to \mathcal{X} \times \mathcal{X}$
est représentable par des espaces algébriques\footnote{Ici, on peut entendre
soit des faisceaux pour la topologie étale dont la diagonale est représentable et qui
admettent un recouvrement étale surjectif par un schéma, soit des espaces algébriques
au sens de
Algebraic Spaces, définition \ref{spaces-definition-algebraic-space}.
En effet, d'après Bootstrap, lemme \ref{bootstrap-lemma-spaces-etale},
il n'y a aucune différence.}, et
\item il existe $U \in \Ob(\Sch_\alpha/S)$
et un $1$-morphisme $(\Sch/U)_\etale \to \mathcal{X}$
qui est surjectif et lisse.
\end{enumerate}
Alors $\mathcal{X}$ est un champ algébrique au sens de
Algebraic Stacks, définition \ref{algebraic-definition-algebraic-stack}.
\end{lemma}

\begin{proof}
Notons que les propriétés (2) et (3) du lemme et les propriétés
(2) et (3) correspondantes de
Algebraic Stacks, définition \ref{algebraic-definition-algebraic-stack},
sont indépendantes de la topologie. En effet, ces propriétés
ne font intervenir que la notion de $2$-produit fibré de catégories fibrées en
groupoïdes, les $1$- et $2$-morphismes de catégories fibrées en groupoïdes, la
notion de $1$-morphisme de catégories fibrées en groupoïdes représentable
par des espaces algébriques et ce que signifie, pour un tel $1$-morphisme, être
surjectif et lisse.
Il suffit donc de montrer qu'un champ en groupoïdes $\mathcal{X}$ pour la topologie
étale qui possède les propriétés (2) et (3) est aussi un champ en groupoïdes pour la topologie fppf.

\medskip\noindent
D'après (2), soit $R$ un espace algébrique représentant
$$
(\Sch_\alpha/U) \times_\mathcal{X} (\Sch_\alpha/U)
$$
D'après (3), les projections $s, t : R \to U$ sont lisses. Exactement comme dans la preuve de
Algebraic Stacks, lemme \ref{algebraic-lemma-map-space-into-stack},
il existe un groupoïde en espaces $(U, R, s, t, c)$ et un $1$-morphisme
canonique pleinement fidèle $[U/R]_\etale \to \mathcal{X}$,
où $[U/R]_\etale$ est le champ associé pour la topologie étale au préfaisceau
en groupoïdes
$$
T \longmapsto (U(T), R(T), s(T), t(T), c(T))
$$
Affirmation : si $V \to T$ est un morphisme lisse et surjectif d'un espace algébrique
$V$ vers un schéma $T$, alors il existe un recouvrement étale $\{T_i \to T\}$
plus fin que le recouvrement $\{V \to T\}$. Cela résulte de
More on Morphisms, lemme \ref{more-morphisms-lemma-etale-dominates-smooth}
ou du résultat plus général
Sheaves on Stacks, lemme
\ref{stacks-sheaves-lemma-surjective-flat-locally-finite-presentation}.
En utilisant l'affirmation et en raisonnant exactement comme dans
Algebraic Stacks, lemme \ref{algebraic-lemma-stack-presentation},
on voit que $[U/R]_\etale \to \mathcal{X}$ est une
équivalence.

\medskip\noindent
Notons ensuite $[U/R]$ le champ quotient pour la topologie fppf,
qui est un champ algébrique d'après
Algebraic Stacks, théorème
\ref{algebraic-theorem-smooth-groupoid-gives-algebraic-stack}.
On a donc des $1$-morphismes
$$
U \to [U/R]_\etale \to [U/R].
$$
Les morphismes $U \to [U/R]_\etale \cong \mathcal{X}$ et
$U \to [U/R]$ sont tous deux surjectifs et lisses (le premier par hypothèse
et le second d'après le théorème) et, dans les deux cas, les
produits fibrés $U \times_\mathcal{X} U$ et $U \times_{[U/R]} U$
sont représentés par $R$. Par conséquent, le $1$-morphisme
$[U/R]_\etale \to [U/R]$ est pleinement fidèle (puisque les morphismes
dans les champs quotients sont donnés par des morphismes vers $R$, voir
Groupoids in Spaces, section
\ref{spaces-groupoids-section-explicit-quotient-stacks}).

\medskip\noindent
Enfin, pour tout schéma $T$ et tout morphisme $t : T \to [U/R]$, le produit fibré
$V = T \times_{[U/R]} U$ est un espace algébrique surjectif et lisse sur $T$.
D'après l'affirmation ci-dessus, il existe un recouvrement étale $\{T_i \to T\}_{i \in I}$
et des morphismes $T_i \to V$ au-dessus de $T$. Cela prouve que l'objet
$t$ de $[U/R]$ au-dessus de $T$ provient de $U$ localement pour la topologie étale. On en conclut que
$[U/R]_\etale \to [U/R]$ est une équivalence de champs en
groupoïdes sur $(\Sch/S)_\etale$ d'après
Stacks, lemme \ref{stacks-lemma-characterize-essentially-surjective-when-ff}.
Cela achève la preuve.
\end{proof}










\begin{multicols}{2}[\section{Autres chapitres}]
\noindent
Préliminaires
\begin{enumerate}
\item \hyperref[introduction-section-phantom]{Introduction}
\item \hyperref[conventions-section-phantom]{Conventions}
\item \hyperref[sets-section-phantom]{Théorie des ensembles}
\item \hyperref[categories-section-phantom]{Catégories}
\item \hyperref[topology-section-phantom]{Topologie}
\item \hyperref[sheaves-section-phantom]{Faisceaux sur les espaces}
\item \hyperref[sites-section-phantom]{Sites et faisceaux}
\item \hyperref[stacks-section-phantom]{Champs}
\item \hyperref[fields-section-phantom]{Corps}
\item \hyperref[algebra-section-phantom]{Algèbre commutative}
\item \hyperref[brauer-section-phantom]{Groupes de Brauer}
\item \hyperref[homology-section-phantom]{Algèbre homologique}
\item \hyperref[derived-section-phantom]{Catégories dérivées}
\item \hyperref[simplicial-section-phantom]{Méthodes simpliciales}
\item \hyperref[more-algebra-section-phantom]{Compléments d'algèbre}
\item \hyperref[smoothing-section-phantom]{Lissage des morphismes d'anneaux}
\item \hyperref[modules-section-phantom]{Faisceaux de modules}
\item \hyperref[sites-modules-section-phantom]{Modules sur les sites}
\item \hyperref[injectives-section-phantom]{Objets injectifs}
\item \hyperref[cohomology-section-phantom]{Cohomologie des faisceaux}
\item \hyperref[sites-cohomology-section-phantom]{Cohomologie sur les sites}
\item \hyperref[dga-section-phantom]{Algèbre différentielle graduée}
\item \hyperref[dpa-section-phantom]{Algèbre à puissances divisées}
\item \hyperref[sdga-section-phantom]{Faisceaux différentiels gradués}
\item \hyperref[hypercovering-section-phantom]{Hyperrecouvrements}
\end{enumerate}
Schémas
\begin{enumerate}
\setcounter{enumi}{25}
\item \hyperref[schemes-section-phantom]{Schémas}
\item \hyperref[constructions-section-phantom]{Constructions de schémas}
\item \hyperref[properties-section-phantom]{Propriétés des schémas}
\item \hyperref[morphisms-section-phantom]{Morphismes de schémas}
\item \hyperref[coherent-section-phantom]{Cohomologie des schémas}
\item \hyperref[divisors-section-phantom]{Diviseurs}
\item \hyperref[limits-section-phantom]{Limites de schémas}
\item \hyperref[varieties-section-phantom]{Variétés}
\item \hyperref[topologies-section-phantom]{Topologies sur les schémas}
\item \hyperref[descent-section-phantom]{Descente}
\item \hyperref[perfect-section-phantom]{Catégories dérivées de schémas}
\item \hyperref[more-morphisms-section-phantom]{Compléments sur les morphismes}
\item \hyperref[flat-section-phantom]{Compléments sur la platitude}
\item \hyperref[groupoids-section-phantom]{Schémas en groupoïdes}
\item \hyperref[more-groupoids-section-phantom]{Compléments sur les schémas en groupoïdes}
\item \hyperref[etale-section-phantom]{Morphismes étales de schémas}
\end{enumerate}
Thèmes de théorie des schémas
\begin{enumerate}
\setcounter{enumi}{41}
\item \hyperref[chow-section-phantom]{Homologie de Chow}
\item \hyperref[intersection-section-phantom]{Théorie de l'intersection}
\item \hyperref[pic-section-phantom]{Schémas de Picard des courbes}
\item \hyperref[weil-section-phantom]{Théories cohomologiques de Weil}
\item \hyperref[adequate-section-phantom]{Modules adéquats}
\item \hyperref[dualizing-section-phantom]{Complexes dualisants}
\item \hyperref[duality-section-phantom]{Dualité pour les schémas}
\item \hyperref[discriminant-section-phantom]{Discriminants et différentes}
\item \hyperref[derham-section-phantom]{Cohomologie de de Rham}
\item \hyperref[local-cohomology-section-phantom]{Cohomologie locale}
\item \hyperref[algebraization-section-phantom]{Géométrie algébrique et formelle}
\item \hyperref[curves-section-phantom]{Courbes algébriques}
\item \hyperref[resolve-section-phantom]{Résolution des surfaces}
\item \hyperref[models-section-phantom]{Réduction semi-stable}
\item \hyperref[functors-section-phantom]{Foncteurs et morphismes}
\item \hyperref[equiv-section-phantom]{Catégories dérivées de variétés}
\item \hyperref[pione-section-phantom]{Groupes fondamentaux de schémas}
\item \hyperref[etale-cohomology-section-phantom]{Cohomologie étale}
\item \hyperref[crystalline-section-phantom]{Cohomologie cristalline}
\item \hyperref[proetale-section-phantom]{Cohomologie pro-étale}
\item \hyperref[relative-cycles-section-phantom]{Cycles relatifs}
\item \hyperref[more-etale-section-phantom]{Compléments de cohomologie étale}
\item \hyperref[trace-section-phantom]{La formule des traces}
\end{enumerate}
Espaces algébriques
\begin{enumerate}
\setcounter{enumi}{64}
\item \hyperref[spaces-section-phantom]{Espaces algébriques}
\item \hyperref[spaces-properties-section-phantom]{Propriétés des espaces algébriques}
\item \hyperref[spaces-morphisms-section-phantom]{Morphismes d'espaces algébriques}
\item \hyperref[decent-spaces-section-phantom]{Espaces algébriques décents}
\item \hyperref[spaces-cohomology-section-phantom]{Cohomologie des espaces algébriques}
\item \hyperref[spaces-limits-section-phantom]{Limites d'espaces algébriques}
\item \hyperref[spaces-divisors-section-phantom]{Diviseurs sur les espaces algébriques}
\item \hyperref[spaces-over-fields-section-phantom]{Espaces algébriques sur des corps}
\item \hyperref[spaces-topologies-section-phantom]{Topologies sur les espaces algébriques}
\item \hyperref[spaces-descent-section-phantom]{Descente et espaces algébriques}
\item \hyperref[spaces-perfect-section-phantom]{Catégories dérivées d'espaces}
\item \hyperref[spaces-more-morphisms-section-phantom]{Compléments sur les morphismes d'espaces}
\item \hyperref[spaces-flat-section-phantom]{Platitude sur les espaces algébriques}
\item \hyperref[spaces-groupoids-section-phantom]{Groupoïdes dans les espaces algébriques}
\item \hyperref[spaces-more-groupoids-section-phantom]{Compléments sur les groupoïdes dans les espaces}
\item \hyperref[bootstrap-section-phantom]{Amorçage}
\item \hyperref[spaces-pushouts-section-phantom]{Sommes amalgamées d'espaces algébriques}
\end{enumerate}
Thèmes de géométrie
\begin{enumerate}
\setcounter{enumi}{81}
\item \hyperref[spaces-chow-section-phantom]{Groupes de Chow des espaces}
\item \hyperref[groupoids-quotients-section-phantom]{Quotients de groupoïdes}
\item \hyperref[spaces-more-cohomology-section-phantom]{Compléments sur la cohomologie des espaces}
\item \hyperref[spaces-simplicial-section-phantom]{Espaces simpliciaux}
\item \hyperref[spaces-duality-section-phantom]{Dualité pour les espaces}
\item \hyperref[formal-spaces-section-phantom]{Espaces algébriques formels}
\item \hyperref[restricted-section-phantom]{Algébrisation des espaces formels}
\item \hyperref[spaces-resolve-section-phantom]{Résolution des surfaces revisitée}
\end{enumerate}
Théorie des déformations
\begin{enumerate}
\setcounter{enumi}{89}
\item \hyperref[formal-defos-section-phantom]{Théorie formelle des déformations}
\item \hyperref[defos-section-phantom]{Théorie des déformations}
\item \hyperref[cotangent-section-phantom]{Le complexe cotangent}
\item \hyperref[examples-defos-section-phantom]{Problèmes de déformation}
\end{enumerate}
Champs algébriques
\begin{enumerate}
\setcounter{enumi}{93}
\item \hyperref[algebraic-section-phantom]{Champs algébriques}
\item \hyperref[examples-stacks-section-phantom]{Exemples de champs}
\item \hyperref[stacks-sheaves-section-phantom]{Faisceaux sur les champs algébriques}
\item \hyperref[criteria-section-phantom]{Critères de représentabilité}
\item \hyperref[artin-section-phantom]{Axiomes d'Artin}
\item \hyperref[quot-section-phantom]{Espaces de Quot et de Hilbert}
\item \hyperref[stacks-properties-section-phantom]{Propriétés des champs algébriques}
\item \hyperref[stacks-morphisms-section-phantom]{Morphismes de champs algébriques}
\item \hyperref[stacks-limits-section-phantom]{Limites de champs algébriques}
\item \hyperref[stacks-cohomology-section-phantom]{Cohomologie des champs algébriques}
\item \hyperref[stacks-perfect-section-phantom]{Catégories dérivées de champs}
\item \hyperref[stacks-introduction-section-phantom]{Introduction aux champs algébriques}
\item \hyperref[stacks-more-morphisms-section-phantom]{Compléments sur les morphismes de champs}
\item \hyperref[stacks-geometry-section-phantom]{La géométrie des champs}
\end{enumerate}
Thèmes de théorie des espaces de modules
\begin{enumerate}
\setcounter{enumi}{107}
\item \hyperref[moduli-section-phantom]{Champs de modules}
\item \hyperref[moduli-curves-section-phantom]{Espaces de modules de courbes}
\end{enumerate}
Divers
\begin{enumerate}
\setcounter{enumi}{109}
\item \hyperref[examples-section-phantom]{Exemples}
\item \hyperref[exercises-section-phantom]{Exercices}
\item \hyperref[guide-section-phantom]{Guide bibliographique}
\item \hyperref[desirables-section-phantom]{Desiderata}
\item \hyperref[coding-section-phantom]{Style de programmation}
\item \hyperref[obsolete-section-phantom]{Obsolète}
\item \hyperref[fdl-section-phantom]{Licence de documentation libre GNU}
\item \hyperref[index-section-phantom]{Index généré automatiquement}
\end{enumerate}
\end{multicols}


\bibliography{my}
\bibliographystyle{amsalpha}

\end{document}
